\documentclass[12pt]{article}
\usepackage{amssymb,amsmath,amsthm,secdot,bbm}

\usepackage[
bookmarks=true,
bookmarksnumbered=true,
colorlinks=true, pdfstartview=FitV, linkcolor=blue, citecolor=blue,urlcolor=blue]{hyperref}


\usepackage[shortlabels]{enumitem}
\usepackage{color}

\usepackage[nameinlink,capitalize]{cleveref}


\topmargin = -2cm %
\oddsidemargin =0cm%
\textwidth = 16cm%
\textheight= 24.5cm%

\marginparwidth=57pt

\crefname{Lemma}{Lemma}{Lemmas}
\crefname{Theorem}{Theorem}{Theorems}
\crefrangeformat{proposition}{Propositions~#3#1#4--#5#2#6}

\usepackage{thmtools}

\newtheorem{theorem}{Theorem}
\newtheorem{lemma}[theorem]{Lemma}
\newtheorem{proposition}[theorem]{Proposition}

\newtheorem{corollary}[theorem]{Corollary}

\theoremstyle{definition}
\newtheorem{assumption}[theorem]{assumption}
\newtheorem{remark}[theorem]{Remark}
\newtheorem{definition}[theorem]{Definition}
\newtheorem{example}[theorem]{Example}





\newcommand{\var}{\textrm{var}}

\newcommand*{\Gref}[1]{\hyperref[SDEfunc]{Eq($#1$)}}
\newcommand*{\Gtag}[1]{\tag{Eq($#1$)}}
\DeclareMathOperator{\1}{\mathbbm{1}}%%
\DeclareMathOperator{\I}{\mathbbm{1}}%%
\DeclareMathOperator{\law}{Law}%
\DeclareMathOperator*{\esssup}{ess\,sup}
\DeclareMathOperator*{\re}{Re}

\def\E{\hskip.15ex\mathsf{E}\hskip.10ex}
\def\P{\mathsf{P}}
\def\Var{\mathop{\mbox{\rm Var}}}

\def\eps{\varepsilon}
\def\phi{\varphi}

\renewcommand{\d}{\partial}

\renewcommand{\ge}{\geqslant}
\renewcommand{\le}{\leqslant}

\newcommand{\dis}{\infty}

\newcommand{\nn}{\nonumber}
\newcommand{\wt}{\widetilde}
\newcommand{\wh}{\widehat}

\newcommand{\A}{\mathcal{A}}
\newcommand{\B}{\mathcal{B}}
\newcommand{\C}{\mathcal{C}}
\newcommand{\F}{\mathcal{F}}
\newcommand{\G}{\mathcal{G}}
\newcommand{\N}{\mathbb{N}}
\newcommand{\PP}{\mathcal{P}}
\newcommand{\R}{\mathbb{R}}
\newcommand{\X}{\mathcal{X}}
\newcommand{\Z}{\mathbb{Z}}

\newcommand{\la}{\langle}
\newcommand{\ra}{\rangle}


\newcommand{\Di}[1]{|#1|}

\begin{document}
	
\title{Unique Ergodicity for the stochastic skew heat equation}
	
	\author{
		Oleg Butkovsky%
		\and
		Jonathan C. Mattingly
		\and
		Lorenzo Zambotti
	}
	
	\maketitle

        \begin{center}
           \textbf{Overview of Proof of Main Theorem}  
        \end{center}
We will prove the uniqueness of the invariant measure by showing that
the Markov Semi-group associated with the system is Asymptotic Strong
Feller. We will do this by constructing asymptotic coupling of the two equations
beginning from arbitrary, different initial conditions $u_0$
and $v_0$. We do this by introducing an intermediate process $\wt
v_t$ whose law on path space is equivalent to $v_t$ but which
converges to $u_t$ as $t \rightarrow \infty$.

This is done by in
introducing a drift term of the form $\lambda_t(u_t-\wt v_t)$ to the
equation for $\wt v_t$ so the equation for $\|u_t - \wt v_t\|^2$ has a
contracting term. We essentially  pick $\lambda_t$ to be of comparable
size to $\|u_t - \wt v_t\|^{-\beta}$ for some $\beta>0$. This makes
the equation strongly contractive. Since we will pick $\lambda_t$ to
be a piecewise constant, $\|u_t - \wt v_t\| \rightarrow 0$ only
asymptotically as $t \rightarrow \infty$. To establish the needed
absolue continuity on the infinite time horizon via Girsanov's
theorem, we will have to show that the appropriate infinite-time square
integrality holds.

After the contractiveness and the integrability, the main technical
lemma of the proof is \cref{L26}. This lemma requires some 
stochastic sewing estimates.

        
        \newpage
        \begin{center}
        \textbf{Proof of Main Theorem}  
        \end{center}
Let $D:=[0,1]$, $p$ is the Dirichlet heat kenrel. For $x\in D$, $u_0,v_0\in C_0([0,1])$, $t>0$, we define $V(t,x):=\int_0^t \int_D p_{t-r}(x,y) W(dr,dy)$,
\begin{align}\label{spdeh}
&u_t(x)=P_t u_0(x) +\int_0^t\int_D p_{t-r}(x,y)
        h(u_r(y))\,dydr+V_t(x),\\
&v_t(x)=P_t v_0(x) +\int_0^t\int_D \label{spdehh}
        p_{t-r}(x,y) h(v_r(y))\,dydr+V_t(x),\\
&\wt v_t(x)=P_t
        v_0(x) +\int_0^t\int_D p_{t-r}(x,y) \bigl(h(\wt
        v_r(y))+\lambda_r(u_r(y)-\wt v_r(y))
        \bigr)\,dydr+V_t(x).\label{spdehhh}
\end{align}
Here $\lambda\colon[0,\infty)\to\R_+$ is a deterministic function which is piecewise constant on intervals $[n,n+ 1)$ for $n\in \N$; $h\colon\R\to\R$ is a smooth function with $\|h\|_{C^{-1}}\le 1$.	

If $f\colon\Omega\times [S,T]\times D\to\R$ is a measurable function, then for 
$m\ge1$ we define
\begin{equation}\label{fullsnormsde}
\|f\|_{\C^{0,0}L_m([S,T])}:=\sup_{s\in[S,T]}\sup_{x\in D}\|f_s(x)\|_{L_m(\Omega)}.
\end{equation}	


\bigskip
\noindent\textbf{Function spaces and Gaussian smoothing.} 
For $b\colon\R\to\R$ bounded and measurable and $k\in\N_0$ we use the H\"older spaces $\C^k(\R)$
with norm $\|b\|_{\C^k}=\sum_{j=0}^k\|b^{(j)}\|_{L_\infty}$, and the negative
Besov-type space $\C^{-1}(\R)$ defined through the heat flow:
\begin{equation}\label{eq:negnorm}
\|b\|_{\C^{-1}}:=\sup_{\theta>0}\theta^{1/2}\,\|G_\theta b\|_{L_\infty(\R)},
\qquad
(G_\theta b)(c):=\int_\R g_\theta(c-w)\,b(w)\,dw,\quad 
g_\theta(w):=\tfrac1{\sqrt{2\pi\theta}}e^{-w^2/(2\theta)},
\end{equation}
where $G_\theta$ is the one-dimensional Gaussian semigroup with variance $\theta$ and
$G_0:=\mathrm{Id}$. We record the following elementary smoothing estimate, the
quantitative form of the heat-flow characterization of $\C^{-1}$; it is the
only analytic input we use about the (singular) drift.

\begin{lemma}\label{L:smoothing}
There is a universal constant $C>0$ such that for every bounded measurable
$b\colon\R\to\R$, every $\theta>0$, every $k\in\N_0$ and every $\alpha\in[0,1]$,
\begin{equation}\label{eq:gauss-bound}
\|G_\theta b\|_{\C^{k+\alpha}(\R)}\le C\,\|b\|_{\C^{-1}}\,\theta^{-(k+1+\alpha)/2}.
\end{equation}
In particular $\|G_\theta b\|_{L_\infty}\le \|b\|_{\C^{-1}}\theta^{-1/2}$,
$\|G_\theta b\|_{\C^1}\le C\|b\|_{\C^{-1}}\theta^{-1}$, and
$\|G_\theta b\|_{\C^{2/3}}\le C\|b\|_{\C^{-1}}\theta^{-5/6}$.
\end{lemma}
\begin{proof}
The bound on $\|G_\theta b\|_{L_\infty}$ is the definition \eqref{eq:negnorm}. For $k\ge1$
write, using the semigroup property $G_\theta=G_{\theta/2}\circ G_{\theta/2}$,
$\partial_c^k(G_\theta b)=(\partial_c^k g_{\theta/2})*\,(G_{\theta/2}b)$, so that
\[
\|\partial_c^k(G_\theta b)\|_{L_\infty}
\le \|\partial_c^k g_{\theta/2}\|_{L_1(\R)}\,\|G_{\theta/2}b\|_{L_\infty}
\le c_k\,(\theta/2)^{-k/2}\cdot \|b\|_{\C^{-1}}(\theta/2)^{-1/2}
\le C\|b\|_{\C^{-1}}\theta^{-(k+1)/2},
\]
because $\|\partial_c^k g_{\theta/2}\|_{L_1(\R)}=c_k(\theta/2)^{-k/2}$ by scaling. Summing over
$j\le k$ gives \eqref{eq:gauss-bound} for $\alpha=0$. For $\alpha\in(0,1)$, the $k$-th derivative
$f:=\partial_c^k(G_\theta b)$ satisfies, for any $c\ne c'$,
\[
\frac{|f(c)-f(c')|}{|c-c'|^\alpha}
\le \min\Bigl\{2\|f\|_{L_\infty}|c-c'|^{-\alpha},\ \|f'\|_{L_\infty}|c-c'|^{1-\alpha}\Bigr\}
\le \|f\|_{L_\infty}^{1-\alpha}\|f'\|_{L_\infty}^{\alpha},
\]
and inserting the already-proved bounds $\|f\|_{L_\infty}\le C\|b\|_{\C^{-1}}\theta^{-(k+1)/2}$ and
$\|f'\|_{L_\infty}\le C\|b\|_{\C^{-1}}\theta^{-(k+2)/2}$ yields
$[f]_{\C^\alpha}\le C\|b\|_{\C^{-1}}\theta^{-(k+1+\alpha)/2}$. This proves \eqref{eq:gauss-bound}.
\end{proof}

\noindent\textbf{Conditional Gaussian identity.}
We first record, as a lemma, the conditional-Gaussian structure of the field
$V(t,x)=\int_0^t\int_D p_{t-r}(x,y)\,W(dr,dy)$ that underlies the entire stochastic
sewing argument. Recall $\sigma^2(s,r,y):=\E\bigl(V(r,y)-\E^sV(r,y)\bigr)^2$.

\begin{lemma}\label{L:condgauss}
Fix $0\le s<r$ and $y\in D$. Then:
\begin{enumerate}[(i)]
\item the random variable $N:=V(r,y)-\E^sV(r,y)$ is centered Gaussian with variance
$\sigma^2(s,r,y)$ and is independent of $\F_s$;
\item for every bounded measurable $b\colon\R\to\R$ and every $\F_s$-measurable random
variable $\xi$,
\begin{equation}\label{eq:condgauss}
\E^s\bigl[b\bigl(V(r,y)+\xi\bigr)\bigr]
=\bigl(G_{\sigma^2(s,r,y)}\,b\bigr)\bigl(\E^sV(r,y)+\xi\bigr)\qquad\text{a.s.}
\end{equation}
\end{enumerate}
\end{lemma}
\begin{proof}
\emph{(i)} By definition $V(r,y)=\int_0^r\int_D p_{r-q}(x,y)|_{x=y}\,W(dq,dz)$ is a Wiener
integral of the deterministic kernel $\1_{[0,r]}(q)\,p_{r-q}(y,z)\in L_2([0,r]\times D)$
against the space--time white noise $W$; hence $(V(r,y))$ together with the whole
field $\{W(f):f\in L_2([0,\infty)\times D)\}$ forms a centered Gaussian family. The
filtration is $\F_s=\sigma\{W(g):g\in L_2,\ \mathrm{supp}_t g\subseteq[0,s]\}$
(completed). Split the kernel into its part on $[0,s]$ and on $(s,r]$:
\[
V(r,y)=\underbrace{\int_0^s\!\!\int_D p_{r-q}(y,z)\,W(dq,dz)}_{=:V^{\le s}}
+\underbrace{\int_s^r\!\!\int_D p_{r-q}(y,z)\,W(dq,dz)}_{=:N}.
\]
Here $V^{\le s}$ is $\F_s$-measurable, while $N$ is a Wiener integral of a kernel
supported on the time interval $(s,r]$. Since white-noise stochastic integrals over
disjoint time intervals are jointly Gaussian and \emph{uncorrelated}, $N$ is
uncorrelated with $W(g)$ for every $g$ supported in $[0,s]$; for jointly Gaussian
variables, uncorrelatedness is equivalent to independence, so $N$ is independent of
$\F_s$. Because $\E^s$ is a linear projection in $L_2$ and $V^{\le s}$ is
$\F_s$-measurable while $\E^sN=\E N=0$ (by independence and $\E N=0$), we get
$\E^sV(r,y)=V^{\le s}$, hence $N=V(r,y)-\E^sV(r,y)$. Finally, by the It\^o isometry,
\[
\E N^2=\int_s^r\!\!\int_D p_{r-q}(y,z)^2\,dz\,dq
=\E\bigl(V(r,y)-\E^sV(r,y)\bigr)^2=\sigma^2(s,r,y),
\]
so $N\sim\mathcal N(0,\sigma^2(s,r,y))$ and is independent of $\F_s$. This proves (i).

\emph{(ii)} Write $V(r,y)+\xi=(\E^sV(r,y)+\xi)+N$. Here $\eta:=\E^sV(r,y)+\xi$ is
$\F_s$-measurable and, by (i), $N$ is independent of $\F_s$ with law
$\mathcal N(0,\sigma^2)$, $\sigma^2:=\sigma^2(s,r,y)$. By the substitution/freezing rule
for conditional expectations (\cite[Ch.~II]{ShiryaevProb}), for bounded measurable $b$,
\[
\E^s\bigl[b(\eta+N)\bigr]=\Psi(\eta)\quad\text{a.s.},\qquad
\Psi(c):=\E\bigl[b(c+N)\bigr]=\int_\R b(c+w)\,g_{\sigma^2}(w)\,dw=(G_{\sigma^2}b)(c),
\]
where the second equality uses that $N\sim\mathcal N(0,\sigma^2)$ has density
$g_{\sigma^2}$ and the last is the definition \eqref{eq:negnorm} of $G_{\sigma^2}$
(for $\sigma^2=0$, i.e. $r=s$, set $G_0=\mathrm{Id}$ and the identity is trivial).
Substituting $\eta=\E^sV(r,y)+\xi$ gives \eqref{eq:condgauss}.
\end{proof}

\noindent Identity \eqref{eq:condgauss} (used with $\xi=\E^s\phi(r,y)$, etc.) converts the
singular drift $b$ evaluated along the solution into the \emph{smoothed} drift
$G_{\sigma^2}b$, whose regularity is controlled by \cref{L:smoothing}; this is the
mechanism underlying the proof of \cref{L26}.


We state the following technical result which we prove later in the text.
\begin{lemma}\label{L:main}
Let $m\ge2$,  $\mu\in(0,\frac16)$, Then there exists a constant $C>0$ which depends only on $m$, such that for any $n\in\N$, $n\le s \le t\le n+1$, $x,z\in D$  we have 
\begin{align}\label{sdeb2}
&\Bigl\|\int_s^t \int_D e^{-\lambda_n(t-r)}p_{t-r}(x,y)(h(u_r(y))-h(\wt v_r(y)))\,dy dr\Bigr\|_{L_m(\Omega)}\le  C\lambda_n^{-\frac18}\bigl(\lambda_n^{-1}+ \|u-\wt v\|_{\C^{0,0}L_m([s,t])}\bigr). \\
  	&\Bigl\|\int_s^t \int_D e^{-\lambda_n(t-r)}(p_{t-r}(x,y)-p_{t-r}(z,y))(h(u_r(y))-h(\wt v_r(y)))\,dy dr\Bigr\|_{L_m(\Omega)}\nn\\
	&\qquad\le 
	C |x-z|^{\mu} \bigl(\lambda_n^{-\frac98}+\|u-\wt v\|_{\C^{0,0}L_m([s,t])}\bigr). \label{sdeb3}
\end{align}
\end{lemma}

Now we proceed to the main part of the proof.
\begin{lemma}\label{l:diffl}
Let $n\in\N$, $t\in[n,n+1)$, $x\in D$. Then  we have 
\begin{align}\label{diffeq}
u_t(x)-\wt v_t(x)&=e^{-\lambda_n(t-n)} P_{t-n}(u_n-\wt v_n)(x)\nn\\
&\phantom{=}+ \int_n^t\int_D e^{-\lambda_n (t- r) } p_{t-r}(x,y)\bigl(h(u_r(y))-h(\wt v_r(y))\bigr)\,dy dr. 
\end{align}
\end{lemma}
\begin{proof}
Let $t\in[n,n+1)$, $x\in D$ and $T>t$. Then 
$$P_{T-t}(u_t-\wt v_t)(x)=P_{T-n}(u_n-\wt v_n)(x)+\int_n^t\int_D p_{T-r}(x,y) \bigl(h(u_r(y))-h(\wt v_r(y))-\lambda_n(u_r(y)-\wt v_r(y))\bigr)\,dydr.
$$ 
Therefore, using the chain rule (in $t$ for fixed $x\in D$) we derive 	
\begin{align*}
e^{\lambda_n t}P_{T-t}(u_t-\wt v_t)(x)&=e^{n \lambda_n } P_{T-n}(u_n-\wt v_n)(x)+
	\lambda_n\int_n^t e^{\lambda_n r } P_{T-r}(u_r-\wt v_r)(x)\,dr\\
		&\phantom{=}+
		\int_n^t\int_D e^{\lambda_n r } p_{T-r}(x,y) \bigl(h(u_r(y))-h(\wt v_r(y))-\lambda_n(u_r(y)-\wt v_r(y))\bigr)\,dydr\\
		&=e^{n \lambda_n } P_{T-n}(u_n-\wt v_n)(x)+ \int_n^t\int_D e^{\lambda r } p_{T-r}(x,y)\bigl(h(u_r(y))-h(\wt v_r(y))\bigr)\,dy dr. 
	\end{align*}
Thus, by letting $T\searrow t$ and using continuity of $u$, $\wt v$ we get \eqref{diffeq}.
\end{proof}
 
\begin{lemma}\label{L:41n}
Let $m\ge2$. There exists a universal constant $C=C(m)$ with the following property:
if $\inf_{n\in\N}\lambda_n\ge C$, then for any $n\in\N$ we have
\begin{equation*}
\|u-\wt v\|_{\C^{0,0}L_m([n,n+1])}\le  2 \sup_{x\in D}\|u_n(x)-v_n(x)\|_{L_m(\Omega)}+C\lambda_n^{-\frac98}. 
\end{equation*}
\end{lemma}

\begin{proof}
  For any $t>0$, $x\in D$ and any measurable function $f\colon\Omega\times D\to\R$,
\begin{equation}\label{ineq}
\|P_t f(x)\|_{L_m(\Omega)}\le 
\int_D p_t(x,y) \|f(y)\|_{L_m(\Omega)}\,dy \le \sup_{y\in D} \|f(y)\|_{L_m(\Omega)}. 
\end{equation}
Using this inequality,   \eqref{diffeq} and taking $s=n$ in \eqref{sdeb2}, we derive for $t\in[n,n+1]$, $x\in D$
\begin{align*}
\|u_t(x)-\wt v_t(x)\|_{L_m(\Omega)}\le \sup_{x\in D}\|u_n(x)-v_n(x)\|_{L_m(\Omega)}+C \lambda_n^{-\frac98}+C \lambda_n^{-\frac18} \|u-\wt v\|_{\C^{0,0}L_m([n,n+1])}
\end{align*}
for $C=C(m)>0$. 
Therefore, taking in the above inequality supremum over $t\in[n,n+1]$, $x\in D$, we get 
$$\|u-\wt v\|_{\C^{0,0}L_m([n,n+1])}\le \sup_{x\in D}\|u_n(x)-v_n(x)\|_{L_m(\Omega)}+C\lambda_n^{-\frac98}+C  \lambda_n^{-\frac18} \|u-\wt v\|_{\C^{0,0}L_m([n,n+1])}.
$$ 
We  suppose that for all $n\in\N$ we have
$
C \lambda^{-\frac18}_n\le\frac1{2}.
$
 Then we get  \eqref{ineq}.
\end{proof}
 
Next, we improve  bound \eqref{ineq}. 
\begin{lemma}\label{l:better} Let $m\ge2$. There exist universal constants $\Gamma_0=\Gamma_0(m)$, $\Gamma=\Gamma(m)\ge\Gamma_0^2$ such that if for any $n\in\N$ we have $\lambda_n>\Gamma_0$ then for any $n\in \N$ we have 
\begin{equation*}
\sup_{x\in D}\|u_{n+1}(x)-\wt v_{n+1}(x)\|_{L_m(\Omega)}\le \frac1{4}\sup_{x\in D}\|u_{n}(x)-\wt v_{n}(x)\|_{L_m(\Omega)}+\Gamma \lambda_n^{-\frac98}. 
	\end{equation*}
\end{lemma}

\begin{proof}
Fix $n\in \N$. It follows  from \eqref{diffeq}, inequality \eqref{ineq} and \cref{L:main}, that for any $x\in D$
\begin{align*}%\label{maineqln}
\|u_{n+1}(x)-\wt v_{n+1}(x)\|_{L^m(\Omega)}&\le e^{-\lambda_n} \sup_{x\in D}\|u_n(x)-\wt v_n(x)\|_{L^m(\Omega)}+C \lambda_n^{-\frac18}(\lambda^{-1}+	 \|u-\wt v\|_{\C^{0,0}L^m([n,n+1])})
\\ & \le e^{-\lambda_n}  \sup_{x\in D}\|u_n(x)-\wt v_n(x)\|_{L^m(\Omega)}+C \lambda_n^{-\frac98}+
\\ & +2C \lambda_n^{-\frac18}\bigl( \sup_{x\in D}\|u_n(x)-v_n(x)\|_{L^m(\Omega)}+C  \lambda_n^{-\frac98}\bigr).
\end{align*}
Then we need that $e^{-\lambda_n}+2C \lambda_n^{-\frac18}<\frac14$ to conclude.
\end{proof}

Finally, we pass on to the analysis on $[0,\infty)$. Fix $m\ge2$. Given $u_0$, $v_0$, we define $\lambda_n$, $n\in\Z_+$, as the smallest real number such that
$
\Gamma \lambda_n^{-\frac98}\le 2^{-n-2}\sup_{x\in D}|u_0(x)-v_0(x)|
$ and  $\lambda_n\ge\Gamma_0$ where $\Gamma=\Gamma(m)$, $\Gamma_0=\Gamma_0(m)$ are from \cref{l:better}.
\begin{lemma}\label{L:27}
There exists $C>0$ such that for any $n\in\N$ we have 
\begin{align}\label{firstres}
&\sup_{x\in D}\|u_{n}(x)-\wt v_{n}(x)\|_{L_m(\Omega)}\le 2^{-n}\|u_0-v_0\|_{\dis};\\
&d_{TV}(\law(v_n),\law(\wt v_n))\le C \|u_0-v_0\|_{\dis}^{\frac19};\label{secondres}\\
&\|\sup_{x\in D} |u_{n}(x)-\wt v_{n}(x)|\,\|_{L_m(\Omega)}\le C 2^{-n}\|u_0-v_0\|_{\dis};\label{thirdres}
\end{align}
\end{lemma}
\begin{proof}
  First, let us show \eqref{firstres}. We proceed by induction. Indeed, if $n=0$, \eqref{firstres} is obviously satisfied. Assume \eqref{firstres} holds for $n\in\N$. Then, by \cref{l:better},
$$\sup_{x\in D}\|u_{n+1}(x)-\wt v_{n+1}(x)\|_{L_m(\Omega)}\le \frac1{4}\sup_{x\in D}\|u_{n}(x)-\wt v_{n}(x)\|_{L_m(\Omega)}+\Gamma \lambda_n^{-\frac98}
\le 2^{-n-1}\|u_0-v_0\|_{\dis}
$$
and thus \eqref{firstres} holds.
To show \eqref{secondres}, we note that if $\|u_0-v_0\|_\infty\ge1$ then \eqref{secondres} is obvious. Otherwise, if $\|u_0-v_0\|_\infty\le1$, recalling that $\Gamma\ge\Gamma_0^2$ by \cref{l:better}, we see that 
$
\Gamma \lambda_n^{-\frac98}= 2^{-n-2}\sup_{x\in D}|u_0(x)-v_0(x)|.
$
Note that thanks to \cref{L:41n} an \eqref{firstres}, we have for any $t\ge0$
\begin{equation}\label{anytbound}
\sup_{x\in D}\|u_t(x)-\wt v_t(x)\|_{L_m(\Omega)}\le  2^{-t+2}\sup_{x\in D}|u_0(x)-v_0(x)|.
\end{equation}
Therefore
\begin{align*}
\int_0^\infty\int_D \lambda_t^2 & \E|u_t(x)-\wt v_t(x)|^2\,dxdt\le \int_0^\infty \lambda_t^2 \sup_{x\in D}   \E|u_t(x)-\wt v_t(x)|^2\,dt\\
&\le   \Gamma^{2} \sup_{x\in D}|u_0(x)-v_0(x)|^{\frac29}\int_0^\infty 2^{-\frac29t+8}\, dt =C\sup_{x\in D}|u_0(x)-v_0(x)|^{\frac29}.
\end{align*}
where $\Gamma$ is from \cref{l:better}. Bound \eqref{secondres} follows now from Pinsker's inequality and the Girsanov's theorem, see, e.g., \cite[Lemma A.1, Theorem A.2]{BKS18}.
Finally, to show \eqref{thirdres}, we use \cref{L:main} and \cref{l:diffl}. Let $\mu\in(0,\frac16)$. We derive for any $x,z\in D$
\begin{align*}
&\|u_{n+1}(x)-\wt v_{n+1}(x)-(u_{n+1}(z)-\wt v_{n+1}(z))\|_{L_m(\Omega)}\le \int_D |p_1(x,y)-p_1(z,y)| \|u_{n}(y)-\wt v_{n}(y)\|_{L_m(\Omega)}\,dy\\
	&\phantom{\le}+\Bigl\|\int_{n}^{n+1} \int_{D}e^{-\lambda_n(n+1-r)} (p_{n+1-r}(x,y)-p_{n+1-r}(z,y))(h(u_r(y))-h(\wt v_r(y)))\,dydr\Bigr\|_{L_m(\Omega)}\\
	&\le C|x-z|^{\mu}\bigl(\|u-\wt v\|_{\C^{0,0} L_m([n,n+1])}+\lambda_{n}^{-\frac98}\bigr)
\end{align*}
for $C=C(m,\mu)$. Taking now $\mu=\frac17$ and recalling  \eqref{anytbound} and the definition of $\lambda_n$, we get for any $x,z\in D$
\begin{equation*}
\|u_{n+1}(x)-\wt v_{n+1}(x)-(u_{n+1}(z)-\wt v_{n+1}(z))\|_{L_m(\Omega)}\le C|x-z|^{\frac17} 2^{-n}\|u_0-v_0\|_{\dis}
\end{equation*}
for $C=C(m)$. Choosing now $m$ large enough (it is sufficient to take $m>7$), and applying the Kolmogorov continuity theorem, we obtain \eqref{thirdres}.	
\end{proof}


\begin{proof}[\textbf{Proof of the main result}]
\paragraph{From the mollified to the limiting drift.}
For $m\in\N$ let $h=h^{(m)}:=g_{1/m}$ be the (smooth, bounded) mollification of the singular
drift used in the problem statement, normalized so that $\|h^{(m)}\|_{\C^{-1}}\le 1$ uniformly
in $m$; let $\PP^{m}_t:=P_t^m$ be the corresponding Markov semigroup on $E=C_0([0,1])$, and let
$\PP_t:=P_t$ be the semigroup of the limiting (singular) equation. By hypothesis,
$\PP^m_t\Phi(x)\to\PP_t\Phi(x)$ for every $\Phi\in C_b(E)$ and every deterministic $x\in E$.
All constants in \cref{L:27} are uniform in $m$, because they depend on $h^{(m)}$ only through
$\|h^{(m)}\|_{\C^{-1}}\le1$. We use this uniformity to transfer the contraction estimates to the
limit semigroup $\PP$.

\paragraph{Asymptotic strong Feller property of $\PP$.}
Fix a globally Lipschitz $\phi\colon E\to\R$ with $\|\phi\|_{L_\infty}\le1$ and
$\|\phi\|_{\mathrm{Lip}}\le1$, and fix $u_0,v_0\in E$. Construct, for each fixed $m$, the
coupling $(u^m,v^m,\wt v^{\,m})$ of \eqref{spdeh}--\eqref{spdehhh} driven by the same noise $V$
with $h=h^{(m)}$ and with the feedback gains $\lambda_n$ chosen as before (these depend only on
$u_0,v_0$ and on the $m$-independent constants $\Gamma,\Gamma_0$). By \cref{L:27},
\begin{align}\label{eq:asfm}
|\PP^m_n\phi(u_0)-\PP^m_n\phi(v_0)|
&=|\E\phi(u^m_n)-\E\phi(v^m_n)|\nn\\
&\le |\E\phi(\wt v^{\,m}_n)-\E\phi(v^m_n)|+|\E\phi(u^m_n)-\E\phi(\wt v^{\,m}_n)|\nn\\
&\le \|\phi\|_{L_\infty}\,d_{TV}\bigl(\law(v^m_n),\law(\wt v^{\,m}_n)\bigr)
   +\|\phi\|_{\mathrm{Lip}}\,\E\|u^m_n-\wt v^{\,m}_n\|^{\frac19}_{\infty}\nn\\
&\le C\,\|u_0-v_0\|_\infty^{\frac19}+C\,2^{-n/9}\|u_0-v_0\|_\infty^{\frac19},
\end{align}
where in the last line we used \eqref{secondres} (Pinsker's inequality plus Girsanov's
theorem, the required infinite-horizon square-integrability being established in the proof of
\eqref{secondres} via \eqref{anytbound}) and \eqref{thirdres}, and the constant $C$ depends only on the fixed integrability exponent used in \cref{L:27} and is \emph{independent of the mollification parameter $m$}. Letting $m\to\infty$
in \eqref{eq:asfm} and using $\PP^m_n\phi\to\PP_n\phi$ pointwise (valid for the bounded
continuous $\phi$ here) we obtain, for the limit semigroup $\PP$,
\begin{equation}\label{eq:asf-limit}
|\PP_n\phi(u_0)-\PP_n\phi(v_0)|\le C\,\|u_0-v_0\|_\infty^{\frac19}+C\,2^{-n/9}\|u_0-v_0\|_\infty^{\frac19},
\qquad n\in\N .
\end{equation}
Inequality \eqref{eq:asf-limit} holds for every $1$-Lipschitz, $1$-bounded $\phi$. Given
$\varepsilon>0$, choosing $n$ with $C2^{-n/9}<\varepsilon$ and then $\delta>0$ with
$C(1+\varepsilon)\delta^{1/9}<\varepsilon$, we get that $\|u_0-v_0\|_\infty<\delta$ implies
$|\PP_n\phi(u_0)-\PP_n\phi(v_0)|<2\varepsilon$ uniformly over the unit ball of Lipschitz
functions. This is exactly the asymptotic strong Feller (ASF) property of $\PP$ at every
point (with the totally bounded gauge $\delta_n\to0$ and a sequence $t_n=n\to\infty$); see
\cite[Definition~3.8 and Proposition~3.12]{HM06}.

\paragraph{Existence of an invariant measure.}
We verify the Krylov--Bogolyubov tightness criterion for $\PP$ by exhibiting a uniform-in-time
\emph{spatial} H\"older bound on the solution. Start the equation \eqref{spdeh} from
$u_0=0$. For $0\le t\le 1$ and $x,z\in D$ we have, by the mild formulation and the same
arguments used to derive \eqref{thirdres},
\begin{equation}\label{eq:spatial-incr}
\|u_t(x)-u_t(z)\|_{L_m(\Omega)}
\le \underbrace{\|V_t(x)-V_t(z)\|_{L_m(\Omega)}}_{(\mathrm I)}
+\underbrace{\Bigl\|\int_0^t\!\!\int_D\bigl(p_{t-r}(x,y)-p_{t-r}(z,y)\bigr)h^{(m)}(u_r(y))\,dy\,dr\Bigr\|_{L_m(\Omega)}}_{(\mathrm{II})} .
\end{equation}
For the Gaussian part $(\mathrm I)$, the It\^o isometry and the standard Dirichlet heat-kernel
estimate $\int_0^t\!\int_D(p_{t-r}(x,y)-p_{t-r}(z,y))^2\,dy\,dr\le C|x-z|$ give, after
Gaussian hypercontractivity, $(\mathrm I)\le C_m|x-z|^{1/2}$ for every $m\ge2$, uniformly in
$t\in[0,1]$. For the drift part $(\mathrm{II})$, applying \cref{c:39}--\cref{c:310} (the
H\"older-in-space estimate \eqref{holder}/\eqref{holderex}) with $\phi=u-V$, $\psi=0$,
$\|h^{(m)}\|_{\C^{-1}}\le1$ and the a priori bound $[u-V]_{\X}\le C(1+\|h^{(m)}\|^2_{\C^{-1}})\le C$
of \cref{l:apriori}, we obtain $(\mathrm{II})\le C|x-z|^{\mu}$ for any $\mu\in(0,\tfrac16)$,
uniformly in $t\in[0,1]$ and in $m$. Choosing $\mu=\tfrac17$ and $m>7$ and invoking the
Kolmogorov continuity theorem in the spatial variable, we get a constant $C$, independent of
$m$ and of $t\in[0,1]$, with $\E\|u_t\|_{\C^{1/9}([0,1])}\le C$. The same estimate run on each
interval $[n,n+1]$, started from the $\F_n$-measurable datum $u_n$ (whose contribution to the
mild formula is the contraction term $P_{t-n}u_n$, which only \emph{improves} spatial
regularity, since $\|P_a f\|_{\C^{1/9}}\le \|f\|_{\C^{1/9}}$), yields the uniform-in-time bound
\begin{equation}\label{eq:holder-unif}
\sup_{t\ge0}\ \E\,\|u_t\|_{\C^{1/9}([0,1])}\le C<\infty .
\end{equation}
(Here we used that $u_t(0)=u_t(1)=0$ by the Dirichlet boundary condition, so the H\"older
seminorm controls the full $C_0([0,1])$-norm: $\|u_t\|_\infty\le \|u_t\|_{\C^{1/9}}$.)
Since the embedding $\C^{1/9}([0,1])\cap C_0([0,1])\hookrightarrow C_0([0,1])$ is compact,
\eqref{eq:holder-unif} and Markov's inequality show that the time-averaged laws
$\bar\mu_T:=\frac1T\int_0^T \law(u_t)\,dt$ ($T\ge1$) are tight on $E=C_0([0,1])$: for the
compact set $K_R=\{w\in C_0([0,1]):\|w\|_{\C^{1/9}}\le R\}$,
\[
\bar\mu_T(E\setminus K_R)\le \frac1T\int_0^T\frac{\E\|u_t\|_{\C^{1/9}}}{R}\,dt\le\frac{C}{R}\xrightarrow[R\to\infty]{}0,
\]
uniformly in $T$. By Prokhorov's theorem $(\bar\mu_T)$ has a weakly convergent subsequence with
limit $\mu_\star$, and by the Krylov--Bogolyubov theorem $\mu_\star$ is invariant for $\PP$,
since $\PP$ is Feller on $E$ (it is the pointwise limit of the Feller semigroups $\PP^m$ and
acts on the path space of the global mild solution granted by the problem statement). Hence
$\PP$ admits at least one invariant probability measure on $E$.

\paragraph{Irreducibility and uniqueness.}
The estimates \eqref{secondres} and \eqref{thirdres} are precisely the two ingredients of the
generalized-coupling criterion of Butkovsky--Wunderlich: they yield that $\PP$ is topologically
irreducible (every nonempty open set is reached with positive probability from every starting
point), arguing exactly as in \cite[proof of Theorem~2.7]{BW20}; indeed \eqref{secondres}
provides a $\law(v_n)$-versus-$\law(\wt v_n)$ closeness in total variation that is small for
nearby initial data, and \eqref{thirdres} provides the pathwise contraction making the two
processes meet asymptotically, which together force the supports of $\law(u_n)$ and $\law(v_n)$
to overlap. Combining the ASF property \eqref{eq:asf-limit} with this irreducibility,
\cite[Corollary~3.17]{HM06} (equivalently \cite[Theorem~2.8]{BW20}) shows that $\PP$ has
\emph{at most one} invariant probability measure on $E$.

\paragraph{Conclusion.}
By the previous two paragraphs, $\PP$ has at least one and at most one invariant probability
measure on $E=C_0([0,1])$. Therefore the stochastic skew heat equation has a unique invariant
probability measure.
\end{proof}

It remains now to prove \cref{L:main}. We will use stochastic sewing.

For $0 \le S < T$, consider the simplices 
$\Delta_{[S,T]}:=\{(s,t)\in[S,T]^2\colon s< t\}$ and $\Delta^3_{[S,T]}:=\{(s,u,t)\in[S,T]^3\colon s< u< t\}$. 
If $A_{\cdot,\cdot}$ is a function $\Delta_{[S,T]}\to\R^d$, then we put 
$\delta A_{s,u,t}:= A_{s,t}-A_{s,u}-A_{u,t},\quad  (s,u,t)\in\Delta^3_{[S,T]}$. 
Let $(\F_t)_{t\in[0,T]}$ be the filtration generated by the Wiener
process defining $V$ and denote by  $\E^t[\cdot]:=\E[\cdot|\F_t]$ the conditional expectation given $\F_t$, $t\in[0,T]$. For $x\in D$, $(s,t)\in\Delta_{[0,1]}$ denote
$
\sigma^2(s,t,x)=\E (V(t,x)-\E^s V(t,x))^2.
$

\begin{lemma}\label{L:fkb} There exists $C>0$ such that for any $s,t\in\Delta_{[0,1]}$, $x\in D$ we have 
$\sigma^2(s,t,x)\ge  C ((1-x)\wedge x \wedge (t-s)^{1/2})$.
\end{lemma}


\begin{proof}
Fix $x\in D$ and assume without loss of generality that $x\le 1-x$, i.e. $x=x\wedge(1-x)$.
By the It\^o isometry for the space-time white noise $W$ and the substitution $\rho=t-r$,
\begin{equation}\label{eq:varform}
\sigma^2(s,t,x)=\E\bigl(V(t,x)-\E^sV(t,x)\bigr)^2
=\int_s^t\!\!\int_D p_{t-r}(x,y)^2\,dy\,dr
=\int_0^{t-s}\Bigl(\int_D p_\rho(x,y)^2\,dy\Bigr)\,d\rho .
\end{equation}
By the symmetry $p_\rho(x,y)=p_\rho(y,x)$ of the Dirichlet heat kernel and the
Chapman--Kolmogorov identity $\int_D p_\rho(x,y)\,p_\rho(y,x)\,dy=p_{2\rho}(x,x)$,
the inner integral equals $p_{2\rho}(x,x)$; hence
\begin{equation}\label{eq:varCK}
\sigma^2(s,t,x)=\int_0^{t-s}p_{2\rho}(x,x)\,d\rho
=\tfrac12\int_0^{2(t-s)}p_a(x,x)\,da
\ \ge\ \tfrac12\int_0^{t-s}p_a(x,x)\,da .
\end{equation}
We bound $p_a(x,x)$ from below near the diagonal. The Dirichlet kernel on $D=[0,1]$ for
$\tfrac12\Delta$ admits the method-of-images representation
\[
p_a(x,x)=\frac1{\sqrt{2\pi a}}\sum_{k\in\Z}\Bigl(e^{-(2k)^2/(2a)}-e^{-(2k+2x)^2/(2a)}\Bigr),
\]
obtained from the whole-line Gaussian kernel by reflecting at the endpoints $0,1$. The
$k=0$ term equals $\frac1{\sqrt{2\pi a}}\bigl(1-e^{-2x^2/a}\bigr)$. If $a\le x^2$ then
$1-e^{-2x^2/a}\ge 1-e^{-2}\ge\tfrac34$. The remaining terms ($k\ne0$) are dominated, for
$0<a\le\frac1{16}$, by
\[
\frac1{\sqrt{2\pi a}}\,2\sum_{k\ge1}e^{-(2k-1)^2/(2a)}
\le \frac1{\sqrt{2\pi a}}\,2\,\frac{e^{-1/(2a)}}{1-e^{-4/(2a)}}\le \frac1{\sqrt{2\pi a}}\cdot\tfrac14,
\]
since $e^{-1/(2a)}\le e^{-8}$ for $a\le\frac1{16}$. Therefore there is a universal $c_0>0$ with
\begin{equation}\label{eq:diaglb}
p_a(x,x)\ \ge\ \frac{c_0}{\sqrt a}\qquad\text{whenever } a\le x^2\wedge\tfrac1{16},
\end{equation}
and one may take $c_0=\tfrac{1}{4\sqrt{2\pi}}$. (A direct numerical evaluation of the
spectral series gives $p_a(x,x)\sqrt a\ge 0.25$ throughout this range, comfortably above the
rigorous constant $c_0=\tfrac1{4\sqrt{2\pi}}\approx0.10$, so \eqref{eq:diaglb} holds with room to spare.) Combining \eqref{eq:varCK} and \eqref{eq:diaglb},
\[
\sigma^2(s,t,x)\ \ge\ \tfrac12\int_0^{(t-s)\wedge x^2\wedge\frac1{16}}\frac{c_0}{\sqrt a}\,da
=c_0\sqrt{(t-s)\wedge x^2\wedge\tfrac1{16}}
\ \ge\ C\bigl((t-s)^{1/2}\wedge x\bigr),
\]
using $\sqrt{x^2\wedge\frac1{16}}\ge \tfrac14 x$ for $x\le\tfrac12$ (indeed $x^2\le\frac14\le\frac1{16}\cdot4$,
so $\sqrt{x^2\wedge\frac1{16}}\ge \tfrac14 x$). By the symmetric roles of $x$ and $1-x$,
$\sigma^2(s,t,x)\ge C\bigl((1-x)\wedge x\wedge (t-s)^{1/2}\bigr)$, as claimed.
\end{proof}
\begin{lemma}\label{L:skb} For $\rho\in(-2,0]$,
  $\mu\in(0,1]$,  there exists $C=C(\rho,\mu)>0$ so that for any $a\in[0,1]$, $\kappa\in(0,1]$, $x,z\in D$ we have 
\begin{align*}
&\int_D p_{\kappa}(x,y) (y\wedge (1-y)\wedge a )^{\rho}\, dy\le C
  a^\rho (1+\frac{a^2}{\kappa});\\
&\int_D |p_{\kappa}(x,y)-p_{\kappa}(z,y)| (y\wedge (1-y)\wedge a )^{\rho}\, dy\le C|x-z|^\mu  \kappa^{-\frac\mu2} a^\rho (1+\frac{a^2}{\kappa}).
\end{align*}
\end{lemma}
\begin{proof}
The bounds follows by a direct calculation 
using the spectral formula for the Dirichlet heat kernel and the fact that for any $\mu>-1$ there exists a constant $C=C(\mu)>0$ such that for any 
$\kappa\in[0,1]$ we have $\sum_{n=1}^\infty  e^{-n^2 \kappa } n^\mu\le C \kappa^{-\frac12-\frac\mu2}.$
\end{proof}

The next corollary follows from \cref{L:fkb} and that \cref{L:skb}
\begin{corollary}\label{c:key}
 Let $\rho\in(-4,0]$, $\mu\in(0,1]$. Then there exists $C=C(\rho,\mu)>0$ such that for any $x,z\in D$, $(s,t)\in\Delta_{[0,1]}$, $\kappa\in[ t-s,1]$, we have 
$\int_D p_{\kappa}(x,y) \sigma(s,t,y)^{\rho}\, dy\le
C(t-s)^{\frac\rho4}$ and
$\int_D |p_{\kappa}(x,y)-p_{\kappa}(z,y)| \sigma(s,t,y)^{\rho}\, dy\le C|x-z|^{\mu}(t-s)^{\frac\rho4-\frac\mu2}$. 
\end{corollary}

Now we can derive some key integral bounds, which eventually imply \cref{L:main}. For a function $\psi\colon\Omega\times[S,T]\times D\to\R$, $\tau\in(0,1]$, $m\ge1$ we consider the following H\"older-type seminorm 
\begin{equation*}
[f]_{\C^{\tau,0}L_m{[S,T]}}:=\sup_{s,t\in \Delta_{[S,T]}}\sup_{x\in D}\frac{\|f(t,x)-\E^s f(t,x)\|_{L_m(\Omega)}}{(t-s)^\tau};\label{norm1}. 
\end{equation*}
For brevity, we denote $\X:=\C^{\frac34,0}L_m([S,T])$.
\begin{lemma}\label{L26}
Let $m\ge2$. Let $w\colon[0,1]\to\R$ be a continuous function, $\mu\in(0,\frac16)$. 
Then there exists a constant $C=C(\mu,m)>0$ such that for any bounded continuous function $b\colon\R\to\R$, any continuous adapted processes  $\phi,\psi\colon\Omega\times[0,1]\times D\to\R$, $x,z\in D$, and any $0\le S\le T\le T_0\le 1$ with  $T-S\le T_0-T$
setting  $\Gamma:=C \|b\|_{\C^{-1}}\|w\|_{L_\infty([S,T])}$,
\begin{equation}\label{cagerm}
\A_{t}^\phi(x):=\int_0^t \int_D w_r \, p_{T_0-r}(x,y) \, b( V_r(y)+\phi_r(y))d r d y,\quad t\in [0,T],
\end{equation}
and setting $\A^{\phi,\psi}_{t}(x):=\A^\phi_t(x)-\A^\psi_t(x)$, $\A^{\phi,\psi}_{s,t}(x):=\A^{\phi,\psi}_{t}(x)-\A^{\phi,\psi}_{s}(x)$ we have
\begin{align}\label{mainres} 
&\|\A^\phi_T(x)-\A^\phi_S(x)\|\le  \Gamma  (T-S)^{\frac34}\bigl(1+[\phi]_{\X}(T-S)^{\frac{1}2}\bigr);\\
&\| \A^{\phi,\psi}_{S,T}(x)\|_{L_m(\Omega)}\le  \Gamma (T-S)^{\frac7{12}}\|\psi-\phi\|^{\frac23}_{\C^{0,0}L_m([S,T])}+ \Gamma(T-S)^{\frac{5}4}\bigl([\phi]_{\X}+[\psi]_{\X}\bigr);\label{mainresdif} \\
&\|\A^{\phi,\psi}_{S,T}(x)-\A^{\phi,\psi}_{S,T}(z)\|_{L_m(\Omega)}\le  \Gamma |x-z|^\mu (T-S)^{\frac7{12}-\frac\mu2}\|\psi-\phi\|^{\frac23}_{\C^{0,0}L_m([S,T])}\nn\\
&\hspace{30ex}+ \Gamma |x-z|^\mu                                                                                                                                               (T-S)^{\frac{5}4-\frac\mu2}\bigl([\phi]_{\X}+[\psi]_{\X}\bigr) \label{holderex}. 
\end{align}
\end{lemma}


\medskip
\noindent\textbf{The stochastic sewing lemma and germ-to-process identification.}
We recall the precise form of the stochastic sewing lemma we use, including its
uniqueness clause, and then explain once and for all how the abstract process it
produces is identified with the genuine integral \eqref{cagerm}.

\begin{theorem}[Stochastic sewing lemma, {\cite[Theorem~2.1]{LeSSL}}]\label{T:ssl}
Let $m\ge2$ and $0\le S\le T\le1$. Let $(A_{s,t})_{(s,t)\in\Delta_{[S,T]}}$ be a family
of $\R$-valued random variables such that $A_{s,t}\in L_m(\Omega)$, $A_{s,t}$ is
$\F_t$-measurable for every $(s,t)$, and the map $(s,t)\mapsto A_{s,t}$ is continuous
in $L_m(\Omega)$. Suppose there are constants $C_1,C_2<\infty$ and
$\varepsilon_1,\varepsilon_2>0$ with
\begin{equation}\label{eq:sslhyp}
\|A_{s,t}\|_{L_m(\Omega)}\le C_1(t-s)^{\frac12+\varepsilon_1},
\qquad
\|\E^s\delta A_{s,u,t}\|_{L_m(\Omega)}\le C_2(t-s)^{1+\varepsilon_2},
\end{equation}
for all $(s,u,t)\in\Delta^3_{[S,T]}$. Then there exists a unique (up to
modification) $\F_t$-adapted process $(\mathcal A_t)_{t\in[S,T]}$ with $\mathcal A_S=0$,
continuous in $L_m(\Omega)$, such that, with $\mathcal A_{s,t}:=\mathcal A_t-\mathcal A_s$,
\begin{equation}\label{eq:sslconc}
\|\mathcal A_{s,t}-A_{s,t}\|_{L_m(\Omega)}\le K_2(t-s)^{1+\varepsilon_2},
\qquad
\|\E^s(\mathcal A_{s,t}-A_{s,t})\|_{L_m(\Omega)}\le K_2(t-s)^{1+\varepsilon_2},
\end{equation}
where $K_1=C(m,\varepsilon_1)\,C_1$ and $K_2=C(m,\varepsilon_2)\,C_2$ (in particular the
$L_m$ remainder is controlled by the \emph{germ-consistency} constant $C_2$ with the
\emph{improved} exponent $1+\varepsilon_2>1$, not merely $\tfrac12+\varepsilon_1$).
Moreover, the
process $\mathcal A$ is characterized by \eqref{eq:sslconc}: if $\widehat{\mathcal A}$ is any
$\F_t$-adapted, $L_m$-continuous process with $\widehat{\mathcal A}_S=0$ satisfying, for some
finite constants and some $\widehat\varepsilon_1>0$, $\widehat\varepsilon_2>0$,
\begin{equation}\label{eq:ssluniq}
\|\widehat{\mathcal A}_{s,t}-A_{s,t}\|_{L_m(\Omega)}\le \widehat K_1(t-s)^{\frac12+\widehat\varepsilon_1},
\qquad
\|\E^s(\widehat{\mathcal A}_{s,t}-A_{s,t})\|_{L_m(\Omega)}\le \widehat K_2(t-s)^{1+\widehat\varepsilon_2},
\end{equation}
then $\widehat{\mathcal A}_t=\mathcal A_t$ a.s. for every $t\in[S,T]$.
\end{theorem}

The last (uniqueness) statement is exactly \cite[Theorem~2.1]{LeSSL}: any two
$L_m$-continuous adapted processes whose increments approximate the same germ in the
sense \eqref{eq:ssluniq} differ by a process whose increments are
$o((t-s))$ both in $L_m$ and after $\E^s$, which forces them to coincide (this is the
uniqueness half of the sewing lemma, see \cite[Theorem~2.1, eqs.~(2.3)--(2.4) and the
following remark]{LeSSL}). We will invoke it through the following identification lemma.

\begin{lemma}[Germ-to-process identification]\label{L:ident}
Let $w,b,\phi,T_0,S,T,x$ be as in \cref{L26}, and define the genuine integral
$\A^\phi_t(x)$ by \eqref{cagerm} and the germ $A^\phi_{s,t}$ by \eqref{aagerm}. Assume
the germ bounds \eqref{claim1} hold (with constant $\Gamma=C\|b\|_{\C^{-1}}\|w\|_{L_\infty}$,
uniform over the mollifications $b_n$ below). Then $(\A^\phi_t(x))_{t\in[S,T]}$ is exactly
the process produced by \cref{T:ssl} for the germ $A^\phi$ on $[S,T]$; in particular it
satisfies \eqref{eq:sslconc}, and consequently
\begin{equation}\label{eq:identbound}
\|\A^\phi_T(x)-\A^\phi_S(x)\|_{L_m(\Omega)}\le \|A^\phi_{S,T}\|_{L_m(\Omega)}+K_2(T-S)^{1+\varepsilon_2}.
\end{equation}
The same statement holds verbatim for the germs/processes used in \eqref{mainresdif}
and \eqref{holderex}.
\end{lemma}
\begin{proof}
Write $\widehat{\mathcal A}^b_t:=\A^\phi_t(x)-\A^\phi_S(x)$ for the genuine process built from a
drift $b$, with increments
$\widehat{\mathcal A}^b_{s,t}=\int_s^t\int_D w_r p_{T_0-r}(x,y)\,b(V_r(y)+\phi_r(y))\,dy\,dr$,
and let $A^{b}_{s,t}$ denote the corresponding germ \eqref{aagerm}.

\emph{Step 1: smooth drifts.} Assume first that $b\in C_b^\infty(\R)$ (bounded with bounded
derivatives of all orders). Then $\widehat{\mathcal A}^b$ is, by classical It\^o/Fubini calculus, an
$\F_t$-adapted process, continuous in $L_m(\Omega)$, with $\widehat{\mathcal A}^b_S=0$
(adaptedness and $L_m$-continuity follow as in the discussion preceding \cref{L:fkb}: the
integrand is $\F_r$-measurable and bounded by $\|w\|_{L_\infty}\|b\|_{L_\infty}p_{T_0-r}(x,y)$,
so $\|\widehat{\mathcal A}^b_{s,t}\|_{L_m}\le\|w\|_{L_\infty}\|b\|_{L_\infty}(t-s)$). We check that
$\widehat{\mathcal A}^b$ obeys the SSL characterization \eqref{eq:ssluniq} for the germ $A^b$.

\emph{$L_m$ remainder.} Using $|b|\le\|b\|_{L_\infty}$ and $\int_s^t\int_D p_{T_0-r}\le(t-s)$,
\begin{equation}\label{eq:ident-Lm}
\|\widehat{\mathcal A}^b_{s,t}-A^b_{s,t}\|_{L_m(\Omega)}\le 2\|w\|_{L_\infty}\|b\|_{L_\infty}(t-s)
=:\widehat K_1\,(t-s)^{\frac12+\frac12},
\end{equation}
the first line of \eqref{eq:ssluniq} with $\widehat\varepsilon_1=\tfrac12>0$.

\emph{$\E^s$ remainder.} By definition $A^b_{s,t}=\int_s^t\int_D w_r p_{T_0-r}(x,y)\,
\E^s b(V_r(y)+\E^s\phi_r(y))\,dy\,dr$, while
$\E^s\widehat{\mathcal A}^b_{s,t}=\int_s^t\int_D w_r p_{T_0-r}(x,y)\,\E^s b(V_r(y)+\phi_r(y))\,dy\,dr$.
Hence
\begin{equation}\label{eq:ident-Es}
\E^s\bigl(\widehat{\mathcal A}^b_{s,t}-A^b_{s,t}\bigr)
=\int_s^t\!\!\int_D w_r p_{T_0-r}(x,y)\,\E^s\Bigl[b(V_r(y)+\phi_r(y))-b(V_r(y)+\E^s\phi_r(y))\Bigr]\,dy\,dr .
\end{equation}
For smooth $b$ we estimate the bracket by the mean value theorem. Freezing $y$ and writing
$\Delta_r:=\phi_r(y)-\E^s\phi_r(y)$ (so $\E^s\Delta_r=0$ and, by the definition of the
seminorm $[\,\cdot\,]_{\X}=[\,\cdot\,]_{\C^{3/4,0}L_m([S,T])}$ recalled above,
$\|\Delta_r\|_{L_m}\le[\phi]_{\X}(r-s)^{3/4}$),
\[
b(V_r+\phi_r)-b(V_r+\E^s\phi_r)=\Bigl(\int_0^1 b'\bigl(V_r+\E^s\phi_r+\theta\Delta_r\bigr)\,d\theta\Bigr)\Delta_r ,
\]
so that, taking absolute values and $\E^s$, since $\|b'\|_{L_\infty}<\infty$ for smooth $b$,
\[
\bigl|\E^s\bigl[b(V_r+\phi_r)-b(V_r+\E^s\phi_r)\bigr]\bigr|\le \|b'\|_{L_\infty}\,\E^s|\Delta_r| .
\]
Therefore, by Minkowski and $\|\E^s|\Delta_r|\|_{L_m}\le\|\Delta_r\|_{L_m}\le[\phi]_\X(r-s)^{3/4}$,
\begin{equation}\label{eq:ident-Es-bound}
\bigl\|\E^s(\widehat{\mathcal A}^b_{s,t}-A^b_{s,t})\bigr\|_{L_m(\Omega)}
\le \|w\|_{L_\infty}\|b'\|_{L_\infty}[\phi]_\X\int_s^t(r-s)^{3/4}\,dr\cdot\!\!\sup_{r}\!\int_D p_{T_0-r}\,dy
\le \|w\|_{L_\infty}\|b'\|_{L_\infty}[\phi]_\X\,(t-s)^{7/4}.
\end{equation}
Since $\tfrac74>1$, \eqref{eq:ident-Es-bound} is the second line of \eqref{eq:ssluniq} with
$\widehat\varepsilon_2=\tfrac34>0$ (the constant depends on $\|b'\|_{L_\infty}$, but this is
harmless: it is only used to \emph{identify} the process, while the size of the process is
governed by \eqref{eq:sslconc} with the $\|b\|_{\C^{-1}}$-constants from \eqref{claim1}).
Thus, for smooth $b$, $\widehat{\mathcal A}^b$ satisfies \eqref{eq:ssluniq} for the germ $A^b$, and by
the uniqueness clause of \cref{T:ssl} it equals the SSL process $\mathcal A^b$ of $A^b$. In
particular $\widehat{\mathcal A}^b=\mathcal A^b$ obeys \eqref{eq:sslconc} with the
$\|b\|_{\C^{-1}}$-constants from \eqref{claim1}, giving \eqref{eq:identbound}.

\emph{Step 2: removing smoothness of $b$.} Let $b$ be merely bounded and continuous, as in
\cref{L26}. Let $b_n:=G_{1/n}b$ be its Gaussian mollifications; then $b_n\in C_b^\infty$,
$b_n\to b$ pointwise and boundedly ($\|b_n\|_{L_\infty}\le\|b\|_{L_\infty}$), and
$\|b_n\|_{\C^{-1}}\le\|b\|_{\C^{-1}}$ (the heat-flow norm \eqref{eq:negnorm} is non-increasing
under $G_{1/n}$ by the semigroup property $G_\theta G_{1/n}=G_{\theta+1/n}$ and
$\theta^{1/2}\le(\theta+1/n)^{1/2}$). Consequently the germ bounds \eqref{claim1} hold for
each $b_n$ \emph{with the same constant} $\Gamma=C\|b\|_{\C^{-1}}\|w\|_{L_\infty}$, and by Step
1 the genuine processes $\widehat{\mathcal A}^{b_n}$ equal the SSL processes $\mathcal A^{b_n}$ and obey
\eqref{eq:sslconc} uniformly in $n$. On one hand, by bounded convergence (the integrands
$w_r p_{T_0-r}(x,y)b_n(V_r+\phi_r)$ are dominated by $\|w\|_{L_\infty}\|b\|_{L_\infty}p_{T_0-r}$
and converge a.e.), $\widehat{\mathcal A}^{b_n}_t\to\widehat{\mathcal A}^{b}_t$ in $L_m(\Omega)$ for each
$t$. On the other hand, the germs converge, $A^{b_n}_{s,t}\to A^{b}_{s,t}$ in $L_m$ by the same
bounded convergence applied to \eqref{aagerm}, and the SSL output depends continuously (in
$L_m$) on the germ via the uniform bounds \eqref{eq:sslconc} (the difference of two SSL
processes for germs $A,A'$ is the SSL process of $A-A'$, whose size is controlled by the germ
sizes of $A-A'$; here $A^{b_n}-A^{b_m}\to0$). Passing to the limit, $\widehat{\mathcal A}^{b}$ is
adapted, $L_m$-continuous, vanishes at $S$, and inherits \eqref{eq:sslconc} for the germ $A^b$
with the $\Gamma$-constants. Hence $\widehat{\mathcal A}^b$ is \emph{the} SSL process of $A^b$ and
\eqref{eq:identbound} holds, as claimed.

\emph{Same argument for \eqref{mainresdif} and \eqref{holderex}.} In those cases the germ is
the $\E^s$-projection of the increment of the genuine integral $\A^{\phi,\psi}$ (resp. its
spatial increment), with $\phi$ replaced by $\E^s\phi$ and $\psi$ by $\E^s\psi$ inside the
bounded smooth drift; the $L_m$ remainder \eqref{eq:ident-Lm} (exponent $1$) and the $\E^s$
remainder \eqref{eq:ident-Es-bound} (exponent $7/4$, using
$\|\Delta_r\|_{L_m}\le([\phi]_\X+[\psi]_\X)(r-s)^{3/4}$) are obtained identically for smooth
drifts, and the mollification of Step 2 transfers the identification to bounded continuous
$b$. Thus in every case the SSL output process equals the genuine integral.
\end{proof}


\begin{proof}
	Fix $(S,T,T_0)\in\Delta_{[0,1]}^3$, $x\in D$. 
We begin with the proof of \eqref{mainres}. 
We use the stochastic sewing lemma (SSL) \cite[Theorem~2.1]{LeSSL} and consider the germ 
\begin{equation}\label{aagerm}
A_{s,t}^\phi:=\E^s\int_s^t \int_D w_r p_{T_0-r}(x,y) b( V(r,y)+\E^s
\phi(r,y))\,dr dy,\quad (s,t)\in\Delta_{[S,T]}.
\end{equation}
We claim that for any $(s,u,t)\in\Delta_{[S,T]}$ we have 
\begin{equation}\label{claim1}
\|A_{s,t}^\phi\|_{L_m(\Omega)}\le \Gamma (t-s)^{\frac34};\quad
\|\E^s\delta A^\phi_{s,u,t}\|_{L_m(\Omega)}\le  \Gamma [\phi]_{\X}(t-s)^{\frac54}.
\end{equation}
The two exponents in \eqref{claim1} are $\tfrac34=\tfrac12+\tfrac14$ and
$\tfrac54=1+\tfrac14$, so the SSL hypotheses \eqref{eq:sslhyp} hold with
$\varepsilon_1=\varepsilon_2=\tfrac14>0$, $C_1=\Gamma$, $C_2=\Gamma[\phi]_{\X}$.
Moreover the germ $A^\phi_{s,t}$ in \eqref{aagerm} is $\F_t$-measurable (it is
$\F_s$-measurable, hence a fortiori $\F_t$-measurable) and continuous in $L_m(\Omega)$
in $(s,t)$ by dominated convergence (the integrand is bounded by
$\|w\|_{L_\infty}\|b\|_{L_\infty}p_{T_0-r}(x,y)$). Thus all hypotheses of \cref{T:ssl}
hold once \eqref{claim1} is verified. By the germ-to-process identification
\cref{L:ident}, the process produced by \cref{T:ssl} coincides with the genuine
integral $t\mapsto\A^\phi_t(x)$, and \eqref{eq:identbound} (with $\varepsilon_2=\tfrac14$,
$K_2=C\,C_2=C\,\Gamma[\phi]_{\X}$) together with the first part of \eqref{claim1}
(with $t=T$, $s=S$) yields the desired bound \eqref{mainres}:
\[
\|\A^\phi_T(x)-\A^\phi_S(x)\|_{L_m}\le\Gamma(T-S)^{\frac34}+C\,\Gamma[\phi]_{\X}(T-S)^{\frac54}
= \Gamma(T-S)^{\frac34}\bigl(1+C[\phi]_{\X}(T-S)^{\frac12}\bigr),
\]
which is \eqref{mainres} (after absorbing $C$ into $\Gamma$); here we used
$(T-S)^{5/4}=(T-S)^{3/4}(T-S)^{1/2}$ and that the SSL remainder constant $K_2$ is
proportional to $C_2=\Gamma[\phi]_{\X}$, cf.\ \cite[Theorem~2.1]{LeSSL}.

We begin with the first part of \eqref{claim1}. We have for any $(s,t)\in\Delta_{[S,T]}$
\begin{align*}
|A_{s,t}^\phi|&\le \int_s^t \int_D w_r  p_{T_0-r}(x,y) |\E^s b( V(r,y)+\E^s \phi(r,y))|\,dr dy\\
&= \int_s^t \int_D w_r p_{T_0-r}(x,y) |G_{\sigma^2(s,r,y)} b(\E^s V(r,y)+\E^s\phi(r,y))|\,dr dy\\
&\le \|w\|_{L_\infty([S,T])} \int_s^t \int_D    p_{T_0-r}(x,y) \|G_{\sigma^2(s,r,y)} b\|_{L_\infty(\R)}\,dr dy\\
&\le C\|w\|_{L_\infty([S,T])}  \|b\|_{\C^{-1}}\int_s^t \int_D p_{T_0-r}(x,y) (\sigma(s,r,y))^{-1}\,dr dy\le \Gamma (t-s)^{\frac34}, 
\end{align*}
 where in the last line we used \cref{c:key} with $ \kappa=T_0-r$, $\rho=-1$. Note that $\kappa\ge  T_0-T\ge T-S\ge t-s$ and thus the conditions of \cref{c:key} hold. This shows the first part of  \eqref{claim1}. 
 
To check the second part of  \eqref{claim1}, we write for $(s,u,t)\in\Delta^3_{[S,T]}$
\begin{align*}
|\E^s\delta A^\phi_{s,u,t}|&\le\E^s  \int_u^t \int_D w_rp_{T_0-r}(x,y) |\E^u b(V(r,y)+\E^s \phi(r,y))-\E^u b(V(r,y)+\E^u \phi(r,y))|\,drdy\\&
\le 
\E^s  \int_u^t \int_D w_rp_{T_0-r}(x,y)  \|G_{\sigma^2(u,r,y)} b\|_{\C^1}|\E^s \phi(r,y))-\E^u \phi(r,y))|\,drdy\\
&\le \Gamma
\E^s  \int_u^t \int_D p_{T_0-r}(x,y)  (\sigma(u,r,y))^{-2}|\E^s \phi(r,y))-\E^u \phi(r,y))|\,drdy. 
\end{align*}
Therefore, 
\begin{equation*}
\|\E^s\delta A^\phi_{s,u,t}\|_{L_m(\Omega)}\le \Gamma[\phi]_{\X}(t-s)^{\frac34}
\int_u^t \int_D p_{T_0-r}(x,y)  (\sigma(u,r,y))^{-2}\,drdy\le \Gamma [\phi]_{\X}(t-u)^{\frac54}, 
\end{equation*}
where in the last inequality we used \cref{c:key} with $\kappa=T_0-r$, $\rho=-2>-4$. We see again  that $\kappa\ge  T_0-T\ge T-S\ge t-s$ and thus the conditions of \cref{c:key} hold. This implies the second part of \eqref{claim1}. 
Thus, all the conditions of the stochastic sewing lemma are satisfied and \eqref{mainres} follows from  \cite[Theorem~2.1]{LeSSL}. 


Equation \eqref{mainresdif} follows again from the SSL. Consider the germ $A_{s,t}:=A_{s,t}^\phi-A_{s,t}^\psi$ and process 
$\A_{t}:=\A_{t}^\phi-\A_{s}^\psi$, for $ (s,t)\in\Delta_{[S,T]}$
where $A^\phi$, $\A^\phi$ were define previously and $A^\psi$, $\A^\psi$ are defined analogously. 
We claim now that for any $(s,u,t)\in\Delta_{[S,T]}$ we have 
\begin{equation}\label{claimn1}
	\|A_{s,t}\|_{L_m(\Omega)}\!\le  \Gamma \|\phi-\psi\|^{\frac23}_{\C^{0,0}L_m([S,T])}(t-s)^{\frac7{12}};\,\,
	\|\E^s\delta A_{s,u,t}\|_{L_m(\Omega)}\!\le  \Gamma ([\phi]_{\X}+[\psi]_{\X})(t-s)^{\frac54}.%\label{claimn2}
\end{equation}
We see again that \eqref{claimn1}  implies that all the conditions of the SSL  are satisfied. We also see that the second part of \eqref{claimn1} immediately follows from the second part of \eqref{claim1}. Thus, it remains only to verify the first part of  \eqref{claimn1}. We argue similarly to the first part of the proof and get for any $(s,t)\in\Delta_{[S,T]}$
\begin{align*}
	|A_{s,t}|&\le \int_s^t \int_D w_r  p_{T_0-r}(x,y) |\E^s b( V(r,y)+\E^s \phi(r,y))-\E^s b( V(r,y)+\E^s \psi(r,y))|\,dr dy\\
	&= \int_s^t \int_D w_r p_{T_0-r}(x,y) |G_{\sigma^2(s,r,y)} b(\E^s V(r,y)+\E^s\phi(r,y))-G_{\sigma^2(s,r,y)} b(\E^s V(r,y)+\E^s\psi(r,y))|\,dr dy\\
	&\le \|w\|_{L_\infty([S,T])} \int_s^t \int_D    p_{T_0-r}(x,y) \|G_{\sigma^2(s,r,y)} b\|_{\C^{\frac23}(\R)}|\E^s\phi(r,y)-\E^s\psi(r,y)|^{\frac23}\,dr dy\\
	&\le \Gamma\int_s^t \int_D p_{T_0-r}(x,y) (\sigma(s,r,y))^{-\frac53}|\E^s\phi(r,y)-\E^s\psi(r,y)|^{\frac23}\,dr dy. 
\end{align*}
Hence, using again \cref{c:key} with $ \kappa=T_0-r$, $\rho=-2$, we derive 
\begin{align*}
\|A_{s,t}\|_{L_m(\Omega)}&\le  \Gamma\int_s^t \int_D p_{T_0-r}(x,y) (\sigma(s,r,y))^{-\frac53}\|\E^s\phi(r,y)-\E^s\psi(r,y)\|^{\frac23}_{L_m(\Omega)}\,dr dy\\
&\le  \Gamma\|\phi-\psi\|^{\frac23}_{\C^{0,0}L_m([S,T])}\int_s^t \int_D p_{T_0-r}(x,y) (\sigma(s,r,y))^{-\frac53}\,dr dy\\
&\le  \Gamma\|\phi-\psi\|^{\frac23}_{\C^{0,0}L_m([S,T])}(t-s)^{\frac7{12}}. 
\end{align*}
This shows the first part of  \eqref{claimn1}. 
The germ exponents in \eqref{claimn1} are $\tfrac7{12}=\tfrac12+\tfrac1{12}$ and
$\tfrac54=1+\tfrac14$, so the SSL hypotheses \eqref{eq:sslhyp} hold with
$\varepsilon_1=\tfrac1{12}>0$ and $\varepsilon_2=\tfrac14>0$. The germ
$A_{s,t}=A^\phi_{s,t}-A^\psi_{s,t}$ is $\F_s$-measurable and $L_m$-continuous as before,
so all hypotheses of \cref{T:ssl} hold. By \cref{L:ident} (applied to the genuine
integral $\A^{\phi,\psi}_t(x)=\A^\phi_t(x)-\A^\psi_t(x)$, whose germ \eqref{aagerm} is
the $\E^s$-projection of the increment with $\phi,\psi$ frozen at $\E^s\phi,\E^s\psi$),
the SSL process coincides with $\A^{\phi,\psi}$, and \eqref{mainresdif} follows from the
first inequality of \eqref{eq:sslconc} (with $s=S$, $t=T$) together with the two parts of
\eqref{claimn1}. 
Finally, we show \eqref{holderex}. We use again the stochastic sewing lemma, which we apply similarly to the germ 
\begin{equation*}
A_{s,t}:=\E^s\int_s^t \int_D( p_{T_0-r}(x,y)-p_{T_0-r}(z,y))w_r (b( V(r,y)+\E^s \phi(r,y))-b( V(r,y)+\E^s\psi(r,y)))\,dr dy, 
\end{equation*}
where $(s,t)\in\Delta_{[S,T]}$
and the process $\A^{\phi,\psi}_{t}(x)$.
Arguing exactly as in the first part of the proof and using \cref{c:key}, we derive for any $(s,u,t)\in\Delta_{[S,T]}$ 
$\|A_{s,t}\|_{L_m(\Omega)}\le  \Gamma |x-z|^\mu
\|\phi-\psi\|^{\frac23}_{\C^{0,0}L_m([S,T])}(t-s)^{\frac7{12}-\frac\mu2}$
and $
\|\E^s\delta A^\phi_{s,u,t}\|_{L_m(\Omega)}\le  \Gamma([\phi]_{\X}+[\psi]_{\X})|x-z|^\mu (t-s)^{\frac54-\frac\mu2}
$%\end{align*}
We now check the exponents. The germ for \eqref{holderex} has $L_m$-size of order
$(t-s)^{\frac7{12}-\frac\mu2}$ and $\E^s\delta$-size of order $(t-s)^{\frac54-\frac\mu2}$.
For the SSL we need $\frac7{12}-\frac\mu2>\frac12$ and $\frac54-\frac\mu2>1$, i.e.
$\mu<\frac16$ and $\mu<\frac12$ respectively; the binding constraint is the first one.
For $\mu\in(0,\frac16)$ \emph{strictly}, both inequalities are strict: indeed
$\frac7{12}-\frac\mu2>\frac12\iff\mu<\frac16$, and at $\mu=\frac16$ one has
$\frac7{12}-\frac1{12}=\frac12$, which is \emph{not} admissible, confirming that the
open interval $\mu\in(0,\frac16)$ is exactly the range for which the germ exponent
strictly exceeds $\frac12$. Hence \eqref{eq:sslhyp} holds with
$\varepsilon_1=\frac7{12}-\frac\mu2-\frac12=\frac1{12}-\frac\mu2>0$ and
$\varepsilon_2=\frac54-\frac\mu2-1=\frac14-\frac\mu2>0$. The germ is again
$\F_s$-measurable and $L_m$-continuous, so \cref{T:ssl} applies. By \cref{L:ident}
(applied to the genuine spatial-increment integral
$\A^{\phi,\psi}_t(x)-\A^{\phi,\psi}_t(z)$, whose increment has $\E^s$-projection equal
to the present germ), the SSL process coincides with that genuine integral, and
\eqref{holderex} follows from the first inequality of \eqref{eq:sslconc}. 
\end{proof}

\begin{corollary}\label{c:39}
Suppose that all the conditions of \cref{L26} are satisfied. Then there exists a constant $C=C(m)>0$ such that bounds \eqref{mainres}, \eqref{mainresdif}, \eqref{holderex} hold for any $0\le S\le T\le 1$, $T_0=T$ (without extra restriction $T-S\le T_0-T$). 
\end{corollary}
\begin{proof}
The result follows immediately from \cref{L26} and the taming singularities lemma \cite[Lemma~2.3]{BFG}.
\end{proof}

\begin{corollary}\label{c:310}
Let $m\ge2$. Let $w\colon[0,1]\to\R$ be a continuous function, $\mu\in(0,\frac16)$. 
Then there exists a constant $C=C(\mu,m)>0$ such that for any bounded continuous function $b\colon\R\to\R$, any continuous adapted processes  $\phi,\psi\colon\Omega\times[0,1]\times D\to\R$, $x,z\in D$, and any $0\le S\le T\le 1$ we have 
\begin{align}%\label{keyb} 
&\Bigl\|\int_S^T \int_D w_r p_{T-r}(x,y)\Bigl(b( V(r,y)+\phi(r,y))- b( V(r,y)+\psi(r,y))\Bigr)\,dr dy\Bigr\|_{L_m(\Omega)}\nn\\
&\quad\le  C \|b\|_{\C^{-1}}\|w\|_{L_\infty([S,T])} (T-S)^{\frac14}\|\psi-\phi\|_{\C^{0,0}L_m([S,T])}\nn\\
&\qquad+ C \|b\|_{\C^{-1}}\|w\|_{L_\infty([S,T])}(T-S)^{\frac54}\bigl(1+[\phi]_{\X}+[\psi]_{\X}\bigr);\\
&\Bigl\|\int_S^T \int_D  w_r(p_{T-r}(x,y)-p_{T-r}(z,y)) (b( V(r,y)+\phi(r,y))-b( V(r,y)+\psi(r,y)))\,dr dy\Bigr\|_{L_m(\Omega)}\nn\\
&\quad\le  C \|b\|_{\C^{-1}}\|w\|_{L_\infty([S,T])} |x-z|^\mu \|\psi-\phi\|_{\C^{0,0}L_m([S,T])}\nn\\
&\qquad+ C \|b\|_{\C^{-1}} \|w\|_{L_\infty([S,T])} |x-z|^\mu (T-S)^{\frac{5}4-\frac\mu2}\bigl(1+[\phi]_{\X}+[\psi]_{\X}\bigr) \nn
 \end{align}
\end{corollary}
\begin{proof}
The first inequality follows from  \cref{c:39}, bound \eqref{mainresdif} and the Young inequality ${a_1 a_2\le  a_1^{\frac32} + a_2^3}$ applied to $a_1:=(T-S)^{\frac1{6}}\|\psi-\phi\|_{\C^{0,0}L_m([S,T])}^{\frac23}$, $a_2:=(T-S)^{\frac5{12}}$. 
The second inequality follows in the same way, this time applying the Young inequality to $a_1:=\|\psi-\phi\|_{\C^{0,0}L_m([S,T])}^{\frac23}$, $a_2:=(T-S)^{\frac7{12}-\frac\mu2}$. 
\end{proof}



\begin{corollary}\label{c:finfin}
Let  $m\ge2$, $\mu\in(0,\frac16)$. Then there exists a constant $C=C(\mu,m)>0$ such that for any bounded continuous function $b\colon\R\to\R$, any continuous adapted processes  $\phi,\psi \colon\Omega\times[0,1]\times D\to\R$, $x,z\in D$, $\lambda>1$, and any  $(s,t)\in\Delta_{[0,1]}$ we have
\[
I_{s,t}^\lambda(x):=\int_s^t\int_ D e^{-\lambda(t-r)}p_{t-r}(x,y)\left(b( V_r(y)+\phi_r(y))- b( V_r(y)+\psi_r(y))\right)drd y
\]
\begin{align}\label{keylam}
&\|I_{s,t}^\lambda(x)\|_{L_m(\Omega)}\le 
C \|b\|_{\C^{-1}}                                                                                                                                               \lambda^{-\frac18}\|\phi-\psi\|_{\C^{0,0}L_m([s,t])}+C \|b\|_{\C^{-1}}\lambda^{-\frac{9}8}\bigl(1+[\phi]_{\X}+[\psi]_{\X}\bigr);\\
&\|I_{s,t}^\lambda(x)-I_{s,t}^\lambda(z)\|_{L_m(\Omega)}\le  C \|b\|_{\C^{-1}} |x-z|^\mu                                                                                                                                                      \|\psi-\phi\|_{\C^{0,0}L_m([s,t])}\nn\\
 &\qquad\qquad \qquad \qquad\qquad \qquad+ C \|b\|_{\C^{-1}}  |x-z|^\mu \lambda^{-\frac{9}8}\bigl(1+[\phi]_{\X}+[\psi]_{\X}\bigr). \label{holder}
\end{align}
\end{corollary}
\begin{proof}
Fix $\lambda>1$, let $\eps>0$ be a parameter, which we will choose later. We split the integral in the left hand side of \eqref{keylam} into two: over the time interval $[s,(t-\lambda^{-1+\eps})\vee s]$ and over the interval $[(t-\lambda^{-1+\eps})\vee s, t]$. If $r\le t-\lambda^{-1+\eps}$, then $e^{-\lambda(t-r)}\le e^{-\lambda^\eps}$ and $\|e^{-\lambda (t-\cdot)}\|_{L_\infty([s,(t-\lambda^{-1+\eps})\vee s])}\le e^{-\lambda^\eps}$. Therefore, \cref{c:310} implies 
\begin{align}\label{tb1}
&\Bigl\|\int_s^{(t-\lambda^{-1+\eps})\vee s}\int_ D e^{-\lambda(t-r)}p_{t-r}(x,y)\Bigl(b( V(r,y)+\phi(r,y))- b( V(r,y)+\psi(r,y))\Bigr)\,dr dy\Bigr\|_{L_m(\Omega)}\nn\\
&\quad\le 
C \|b\|_{\C^{-1}} e^{-\lambda^\eps}(1+\|\phi-\psi\|_{\C^{0,0}L_m([s,t])} +[\phi]_{\X}+[\psi]_{\X}\bigr). 
\end{align}
for $C=C(m)$. 
On the other hand, if $r\in [(t-\lambda^{-1+\eps})\vee s, t]$, then $e^{-\lambda(t-r)}\le 1$ and \cref{c:310} implies 
\begin{align}\label{tb2}
&\Bigl\|\int_{(t-\lambda^{-1+\eps})\vee s}^t\int_ D e^{-\lambda(t-r)}p_{t-r}(x,y)\Bigl(b( (V+\phi)(r,y))- b( (V+\psi)(r,y))\Bigr)\,dr dy\Bigr\|_{L_m(\Omega)}\nn\\
&\qquad\le C \|b\|_{\C^{-1}}                                                                                                                                                                     \lambda^{-\frac14(1-\eps)}\|\phi-\psi\|_{\C^{0,0}L_m([s,t])}+C \|b\|_{\C^{-1}} \lambda^{-\frac54(1-\eps)}\bigl(1+[\phi]_{\X}+[\psi]_{\X}\bigr). 
\end{align}
for $C=C(m)$. Choose now $\eps$ such that $\frac54(1-\eps)=\frac{9}8$. Then \eqref{keylam} follows from \eqref{tb1}, \eqref{tb2}, taking into account that for such $\eps$ and some universal $C>0$ we   have $\exp(-\lambda^\eps)\le C \lambda^{-\frac{9}8}$ for all $\lambda\ge1$. 
Bound \eqref{holder} follows from \cref{c:310}  by exactly the same argument. 
\end{proof}	





\begin{lemma}\label{l:apriori}
Let $m\ge2$.  Then there exists a constant $C=C(m)>0$ such that for
any bounded continuous function $b\colon\R\to\R$, a continuous in
space and  $\F_0$-measurable process $z\colon\Omega\times D\to\R$, any  continuous adapted processes $Z,\rho\colon\Omega\times[0,1]\times D\to\R$ satisfying 
\begin{equation*}%\label{zequ}
Z(t,x)=P_t z(x)+ \int_0^t \int_{D} p_{t-r}(x,y)\bigl(b(Z(r,y))+ \rho(r,y)\bigr)\,dydr + V(t,x), t\in[0,1], x\in D, 
\end{equation*}
we have 
$[Z-V]_{\C^{\frac34,0}L_m([0,1])}\le  C (1+\|b\|_{\C^{-1}}^2)(1+\|\rho\|_{ \C^{0,0}L_m([0,1])})$. 
\end{lemma}	

\begin{proof}
Let $\phi_t(x):=Z_t(x)-V_t(x)$, $t\in[0,1]$, $x\in D$, and abbreviate
$R:=\|\rho\|_{\C^{0,0}L_m([0,1])}$. \emph{Throughout the proof we work with the
restricted (local) seminorms}
\[
[\phi]_{\X,[a,c]}:=\sup_{(s',t')\in\Delta_{[a,c]}}\sup_{x\in D}
\frac{\|(\phi_{t'}-\E^{s'}\phi_{t'})(x)\|_{L_m(\Omega)}}{(t'-s')^{3/4}},
\qquad [a,c]\subseteq[0,1],
\]
so that the target quantity is $[\phi]_{\X,[0,1]}=[Z-V]_{\C^{\frac34,0}L_m([0,1])}$.
These seminorms are nondecreasing in the interval: $[a,c]\subseteq[a',c']$ implies
$[\phi]_{\X,[a,c]}\le[\phi]_{\X,[a',c']}$.

\smallskip
\emph{Step 1: per-pair bound.}
Fix $(s,t)\in\Delta_{[0,1]}$ and $x\in D$. Since $\E^s\phi_t$ is, up to the universal
constant $C=C(m)$, the best $L_m(\Omega)$-approximation of $\phi_t$ by an
$\F_s$-measurable random variable, and $P_{t-s}\phi_s$ is $\F_s$-measurable, the
projection property of the conditional expectation in $L_m(\Omega)$
(\cite[Lemma~4.2(i)]{BDGLevy}) gives
\begin{equation}\label{z1step} 	
\|(\phi_t-\E^s \phi_t)(x)\|_{L_m(\Omega)}\le C \|(\phi_t-P_{t-s}\phi_s)(x)\|_{L_m(\Omega)} .
\end{equation}
By the mild equation, $\phi_t-P_{t-s}\phi_s=\int_s^t\int_D p_{t-r}(x,y)\bigl(b(Z_r(y))+\rho_r(y)\bigr)\,dy\,dr$,
so by Minkowski's integral inequality and $\int_Dp_{t-r}(x,y)\,dy\le1$,
\begin{equation}\label{z1bis}
\|(\phi_t-P_{t-s}\phi_s)(x)\|_{L_m(\Omega)}
\le \Bigl\|\int_s^t\!\!\int_D p_{t-r}(x,y) b(Z_r(y))\,dydr\Bigr\|_{L_m(\Omega)}+(t-s)\,R .
\end{equation}
To estimate the first term we apply \cref{c:39} with $w\equiv1$, $T_0=t$, $\psi\equiv0$
on the interval $[s,t]$, recalling $Z_r=V_r+\phi_r$ so that $b(Z_r)=b(V_r+\phi_r)$.
Bound \eqref{mainres} then yields
\begin{equation}\label{z2step}
\Bigl\|\int_s^t\!\!\int_D p_{t-r}(x,y) b(Z_r(y))\,dydr\Bigr\|_{L_m(\Omega)}
\le C \|b\|_{\C^{-1}} (t-s)^{\frac34}\bigl(1+[\phi]_{\X,[s,t]}(t-s)^{\frac12}\bigr),
\end{equation}
for $C=C(m)$, where it is essential that the seminorm appearing on the right is the
\emph{local} one $[\phi]_{\X,[s,t]}$: indeed \cref{c:39}/\cref{L26} produce the germ
bounds \eqref{claim1} on the interval $[s,t]$, whose germ-consistency constant is
$C_2=\Gamma[\phi]_{\X,[s,t]}$. Combining \eqref{z1step}--\eqref{z2step} and dividing by
$(t-s)^{3/4}$ gives, for every $(s,t)\in\Delta_{[0,1]}$ and $x\in D$,
\begin{equation}\label{perpair}
\frac{\|(\phi_t-\E^s\phi_t)(x)\|_{L_m(\Omega)}}{(t-s)^{3/4}}
\le C\|b\|_{\C^{-1}}\bigl(1+(t-s)^{\frac12}[\phi]_{\X,[s,t]}\bigr)+(t-s)^{\frac14}R .
\end{equation}

\smallskip
\emph{Step 2: a priori finiteness of every local seminorm.}
Before any absorption we must know that the local seminorms are finite; otherwise the
term $(t-s)^{1/2}[\phi]_{\X,[s,t]}$ in \eqref{perpair} cannot be moved to the left.
This is the only place where the qualitative hypothesis that $b$ is \emph{bounded} is
used. Set $M:=\|b\|_{L_\infty(\R)}<\infty$. From \eqref{z1step} and the mild equation,
using $\int_Dp_{t-r}(x,y)\,dy\le1$, Minkowski and $|b|\le M$,
\[
\|(\phi_t-\E^s\phi_t)(x)\|_{L_m(\Omega)}
\le C\,(t-s)\bigl(M+R\bigr),
\qquad (s,t)\in\Delta_{[0,1]},\ x\in D .
\]
Dividing by $(t-s)^{3/4}$ and using $t-s\le1$, then taking the supremum over
$(s,t)\in\Delta_{[0,1]}$ and $x\in D$, we get the crude bound
\begin{equation}\label{crude}
[\phi]_{\X,[a,c]}\ \le\ [\phi]_{\X,[0,1]}\ \le\ C\bigl(M+R\bigr)<\infty,
\qquad\text{for every }[a,c]\subseteq[0,1].
\end{equation}
This preliminary estimate depends on $M=\|b\|_{L_\infty}$, which is \emph{not}
controlled by $\|b\|_{\C^{-1}}$ and blows up under mollification; it serves only to
guarantee finiteness so that the self-improving inequality below may be closed. The
final bound depends on $b$ \emph{only} through $\|b\|_{\C^{-1}}$ and is therefore stable
under mollification.

\smallskip
\emph{Step 3: self-improving bound on short intervals.}
Fix an interval $J=[a,c]\subseteq[0,1]$. For every sub-pair
$(s,t)\in\Delta_{J}$ we have $[\phi]_{\X,[s,t]}\le[\phi]_{\X,J}$ and
$(t-s)^{1/2}\le(c-a)^{1/2}$, so \eqref{perpair} gives
\[
\frac{\|(\phi_t-\E^s\phi_t)(x)\|_{L_m(\Omega)}}{(t-s)^{3/4}}
\le C\|b\|_{\C^{-1}}\bigl(1+R\bigr)+C\|b\|_{\C^{-1}}(c-a)^{\frac12}\,[\phi]_{\X,J}.
\]
Taking the supremum over $(s,t)\in\Delta_J$ and $x\in D$ yields the closed inequality
\begin{equation}\label{selfimp}
[\phi]_{\X,J}\le C\|b\|_{\C^{-1}}(1+R)+C\|b\|_{\C^{-1}}(c-a)^{\frac12}\,[\phi]_{\X,J},
\qquad J=[a,c]\subseteq[0,1].
\end{equation}
Crucially, the seminorm on both sides of \eqref{selfimp} is over the \emph{same}
interval $J$, because in \eqref{perpair} the local seminorm $[\phi]_{\X,[s,t]}$ is
controlled by $[\phi]_{\X,J}$ with the \emph{matching} length $(t-s)\le c-a$. Choose
\begin{equation}\label{ellchoice}
\ell:=\bigl(2C\|b\|_{\C^{-1}}\bigr)^{-2},\qquad\text{so that }C\|b\|_{\C^{-1}}\ell^{\frac12}=\tfrac12 .
\end{equation}
Then for any interval $J$ with $|J|=c-a\le\ell$, the coefficient
$C\|b\|_{\C^{-1}}(c-a)^{1/2}\le\frac12$, and since $[\phi]_{\X,J}<\infty$ by
\eqref{crude}, we may absorb it into the left-hand side of \eqref{selfimp}:
\begin{equation}\label{shortbound}
[\phi]_{\X,J}\le 2C\|b\|_{\C^{-1}}(1+R)=:K,
\qquad\text{for every interval }J\subseteq[0,1]\text{ with }|J|\le\ell .
\end{equation}
Note $K\ell^{1/2}=2C\|b\|_{\C^{-1}}(1+R)\cdot(2C\|b\|_{\C^{-1}})^{-1}=1+R$.

\smallskip
\emph{Step 4: from short intervals to the full interval.}
If $\ell\ge1$ then $J=[0,1]$ is already short and \eqref{shortbound} \emph{is} the
conclusion (with $K=2C\|b\|_{\C^{-1}}(1+R)\le C(1+\|b\|_{\C^{-1}}^2)(1+R)$). Suppose
$\ell<1$. Fix an arbitrary pair $(s,t)\in\Delta_{[0,1]}$ and $x\in D$; we bound
$\|(\phi_t-\E^s\phi_t)(x)\|_{L_m}/(t-s)^{3/4}$ by chaining over short sub-intervals.
Let $N:=\lceil (t-s)/\ell\rceil\le\lceil1/\ell\rceil$ and set
$t_k:=s+k(t-s)/N$, $k=0,\dots,N$, so that $d:=t_{k+1}-t_k=(t-s)/N\le\ell$. The
\emph{deterministic} heat semigroup telescopes exactly:
\begin{equation}\label{Ptel}
\phi_t-P_{t-s}\phi_s
=\sum_{k=0}^{N-1} P_{t-t_{k+1}}\bigl(\phi_{t_{k+1}}-P_{t_{k+1}-t_k}\phi_{t_k}\bigr),
\end{equation}
because $P_{t-t_{k+1}}\phi_{t_{k+1}}-P_{t-t_{k+1}}P_{t_{k+1}-t_k}\phi_{t_k}
=P_{t-t_{k+1}}\phi_{t_{k+1}}-P_{t-t_k}\phi_{t_k}$ telescopes to
$\phi_t-P_{t-s}\phi_s$. Each summand is the short-interval mild increment
$\phi_{t_{k+1}}-P_{t_{k+1}-t_k}\phi_{t_k}=\int_{t_k}^{t_{k+1}}\!\int_D
p_{t_{k+1}-r}(\cdot,y)(b(Z_r(y))+\rho_r(y))\,dy\,dr$, whose $L_m$-norm is bounded, by
\eqref{z1bis}--\eqref{z2step} applied on $[t_k,t_{k+1}]$ together with the short-interval
bound \eqref{shortbound} (recall $[t_k,t_{k+1}]$ has length $d\le\ell$, so
$[\phi]_{\X,[t_k,t_{k+1}]}\le K$), by
\[
\bigl\|\phi_{t_{k+1}}-P_{t_{k+1}-t_k}\phi_{t_k}\bigr\|_{L_m(\Omega)}
\le C\|b\|_{\C^{-1}}d^{\frac34}\bigl(1+K d^{\frac12}\bigr)+d\,R
\le C\|b\|_{\C^{-1}}d^{\frac34}\bigl(1+K\ell^{\frac12}\bigr)+d^{\frac34}R ,
\]
where we used $d\le\ell\le1$ and $d\le d^{3/4}$. Since $K\ell^{1/2}=1+R$, the
right-hand side is at most $B\,d^{3/4}$ with $B:=C\|b\|_{\C^{-1}}(2+R)+R\le
C(1+\|b\|_{\C^{-1}})(1+R)$. Using the $L_m$-contractivity of $P_a$
(inequality \eqref{ineq}) on each summand of \eqref{Ptel} and the triangle inequality,
\[
\|(\phi_t-P_{t-s}\phi_s)(x)\|_{L_m(\Omega)}
\le \sum_{k=0}^{N-1}\bigl\|\phi_{t_{k+1}}-P_{t_{k+1}-t_k}\phi_{t_k}\bigr\|_{L_m(\Omega)}
\le N\,B\,d^{\frac34}
= B\,N^{\frac14}(t-s)^{\frac34},
\]
since $N d^{3/4}=N\cdot\bigl((t-s)/N\bigr)^{3/4}=N^{1/4}(t-s)^{3/4}$. Combining with the
projection bound \eqref{z1step} and dividing by $(t-s)^{3/4}$, then taking the supremum
over $(s,t)\in\Delta_{[0,1]}$ and $x\in D$,
\[
[\phi]_{\X,[0,1]}\le C\,B\,N^{\frac14}.
\]
Finally $N\le\lceil1/\ell\rceil\le 1/\ell+1=4C^2\|b\|_{\C^{-1}}^2+1$, so
$N^{1/4}\le 1+(4C^2)^{1/4}\|b\|_{\C^{-1}}^{1/2}\le 1+C\|b\|_{\C^{-1}}$ (using
$\|b\|_{\C^{-1}}^{1/2}\le 1+\|b\|_{\C^{-1}}$). Therefore
\[
[\phi]_{\X,[0,1]}\le C\,(1+\|b\|_{\C^{-1}})(1+R)\cdot(1+\|b\|_{\C^{-1}})
= C\,(1+\|b\|_{\C^{-1}})^2(1+R)\le C\,(1+\|b\|_{\C^{-1}}^2)(1+R),
\]
which is the claimed bound. (The constant $C$ depends only on $m$.)
\end{proof}


\begin{corollary}\label{c:hmain}
Let $u$ be a solution to \eqref{spdeh} and $\wt v$ be a solution to 
\eqref{spdehhh}. Let $m\ge2$, $\mu\in(0,\frac16)$. Then there exists 
$C=C(m,\mu)>0$ such that for any $n\in\N$ and any $n\le s\le t\le n+1$ we have, with
all seminorms understood as the local seminorms of \cref{l:apriori} over $[s,t]$,
\begin{align*}
&[u-V]_{\C^{\frac34,0}L_m([s,t])}\le  C (1+\|h\|_{\C^{-1}}^2);\\
&[\wt v-V]_{\C^{\frac34,0}L_m([s,t])}\le  C (1+\|h\|_{\C^{-1}}^2)\bigl(1+\lambda_n\|u-\wt v\|_{\C^{0,0} L_m([s,t])}\bigr).
\end{align*}
In particular, taking $[s,t]=[n,n+1]$ gives the unit-interval bounds.
\end{corollary}
\begin{proof}
Fix $n\le s\le t\le n+1$. We apply \cref{l:apriori} on the interval $[s,t]$ (after the
time shift $r\mapsto r-s$, which maps $[s,t]$ into $[0,1]$ and under which the entire
proof of \cref{l:apriori} is invariant), with $b=h$, datum
$z=u_s$ (resp. $z=\wt v_s$) at the initial time $s$ (which is $\F_s$-measurable, playing
the role of the $\F_0$-measurable datum after the shift), and the field $V$ replaced by
the $\F_s$-shifted Gaussian field $V(s+\cdot,\cdot)-P_\cdot\,V(s,\cdot)$; the adaptedness
of all processes and the conditional-Gaussian structure are preserved under this shift,
and $\lambda$ is constant $=\lambda_n$ on $[s,t)\subseteq[n,n+1)$. Since the drift $h=h^{(m)}=g_{1/m}$ used throughout the quantitative argument is a
bounded, smooth function (a Gaussian density, so $M:=\|h\|_{L_\infty}<\infty$),
\emph{Step~2} of the proof of \cref{l:apriori} applies and guarantees the a priori
finiteness $[\,\cdot\,]_{\X,[s,t]}<\infty$ that is needed before the self-improving
inequality is closed; the resulting bound depends on $h$ only through
$\|h\|_{\C^{-1}}\le1$ and is therefore uniform in $m$.

For $u$, comparing \eqref{spdeh} with the equation in \cref{l:apriori} we read
off $Z=u$ and $\rho\equiv0$, whence
$[u-V]_{\C^{\frac34,0}L_m([s,t])}\le C(1+\|h\|_{\C^{-1}}^2)$.

For $\wt v$, comparing \eqref{spdehhh} with the equation in \cref{l:apriori} we
read off $Z=\wt v$ and the feedback drift
$\rho(r,y)=\lambda_n\bigl(u_r(y)-\wt v_r(y)\bigr)$ (recall $\lambda$ is constant
$=\lambda_n$ on $[n,n+1)$). Hence
$\|\rho\|_{\C^{0,0}L_m([s,t])}=\lambda_n\|u-\wt v\|_{\C^{0,0}L_m([s,t])}$ and
\cref{l:apriori} gives
\[
[\wt v-V]_{\C^{\frac34,0}L_m([s,t])}
\le C(1+\|h\|_{\C^{-1}}^2)\bigl(1+\lambda_n\|u-\wt v\|_{\C^{0,0}L_m([s,t])}\bigr).
\]
This is the linear-in-$\lambda_n$ growth recorded in the statement. We stress that the
a priori finiteness required to run \cref{l:apriori} for $\wt v$ is supplied by
\emph{Step~2} of its proof: the drift $b=h^{(m)}$ is bounded ($M=\|h^{(m)}\|_{L_\infty}<\infty$)
and, by the global well-posedness granted in the problem statement, $u$ and $\wt v$ are
continuous adapted processes, so the feedback $\rho=\lambda_n(u-\wt v)$ has finite
$\C^{0,0}L_m([s,t])$ norm $\lambda_n\|u-\wt v\|_{\C^{0,0}L_m([s,t])}<\infty$; hence the
crude bound \eqref{crude} gives $[\wt v-V]_{\X,[s,t]}<\infty$ \emph{before} any absorption,
and the displayed estimate is the output of the closed self-improving inequality. The linear
$\lambda_n$ factor is genuine and is cancelled downstream in \cref{L:main} by the
$\lambda_n^{-9/8}$ gain of \cref{c:finfin}.
\end{proof}
\begin{proof}[Proof of \cref{L:main}]
We prove \eqref{sdeb2} by combining \cref{c:finfin} with \cref{c:hmain} and
checking that the linear $\lambda_n$ factor of \cref{c:hmain} is exactly
cancelled by the $\lambda_n^{-9/8}$ gain of \cref{c:finfin}. Fix $n\in\N$ and
$n\le s\le t\le n+1$. Apply \eqref{keylam} of \cref{c:finfin} with $b=h$,
$\lambda=\lambda_n>1$, $\phi=u-V$ and $\psi=\wt v-V$ on $[s,t]\subset[n,n+1]$;
note $h(V_r+\phi_r)=h(u_r)$, $h(V_r+\psi_r)=h(\wt v_r)$ and $\phi-\psi=u-\wt v$,
so $I^{\lambda_n}_{s,t}(x)$ equals the integral on the left of \eqref{sdeb2}.
Thus, writing $U:=\|u-\wt v\|_{\C^{0,0}L_m([s,t])}$ and using $\|h\|_{\C^{-1}}\le1$,
\[
\Bigl\|\!\int_s^t\!\!\int_D e^{-\lambda_n(t-r)}p_{t-r}(x,y)\bigl(h(u_r(y))-h(\wt v_r(y))\bigr)dy\,dr\Bigr\|_{L_m(\Omega)}
\le C\lambda_n^{-\frac18}U+C\lambda_n^{-\frac98}\bigl(1+[\phi]_{\X}+[\psi]_{\X}\bigr).
\]
By \cref{c:hmain} applied on the same interval $[s,t]$ (with $\|h\|_{\C^{-1}}\le1$, so
$1+\|h\|_{\C^{-1}}^2\le2$), the local seminorms over $[s,t]$ satisfy
\[
[\phi]_{\X}=[u-V]_{\C^{\frac34,0}L_m([s,t])}\le C,\qquad
[\psi]_{\X}=[\wt v-V]_{\C^{\frac34,0}L_m([s,t])}\le C\bigl(1+\lambda_n U\bigr),
\]
where $U=\|u-\wt v\|_{\C^{0,0}L_m([s,t])}$, so that the seminorms entering \eqref{keylam}
on $[s,t]$ are exactly the ones bounded by \cref{c:hmain}.
Substituting and expanding (all $\C^{0,0}L_m$ norms taken on $[s,t]\subseteq[n,n+1]$),
\begin{align*}
C\lambda_n^{-\frac18}U+C\lambda_n^{-\frac98}\bigl(1+[\phi]_{\X}+[\psi]_{\X}\bigr)
&\le C\lambda_n^{-\frac18}U+C\lambda_n^{-\frac98}\bigl(1+C+C(1+\lambda_n U)\bigr)\\
&\le C\lambda_n^{-\frac18}U+C\lambda_n^{-\frac98}+\underbrace{C\lambda_n^{-\frac98}\cdot\lambda_n U}_{=\,C\lambda_n^{-1/8}U}\\
&\le C\lambda_n^{-\frac18}U+C\lambda_n^{-\frac98}
= C\lambda_n^{-\frac18}\bigl(\lambda_n^{-1}+U\bigr),
\end{align*}
where in the last identity we crucially used $\lambda_n^{-9/8}=\lambda_n^{-1/8}\cdot\lambda_n^{-1}$.
This is exactly \eqref{sdeb2}; the constant $C$ depends only on $m$ (and on
$\mu$ through \cref{c:finfin}), and the potentially divergent $\lambda_n[\phi]$
contributions have combined into the harmless prefactor $\lambda_n^{-1/8}$.

For \eqref{sdeb3} we argue identically using \eqref{holder} of \cref{c:finfin}
in place of \eqref{keylam}. With the same substitutions, \eqref{holder} yields
\[
\|I^{\lambda_n}_{s,t}(x)-I^{\lambda_n}_{s,t}(z)\|_{L_m(\Omega)}
\le C|x-z|^{\mu}U+C|x-z|^{\mu}\lambda_n^{-\frac98}\bigl(1+[\phi]_{\X}+[\psi]_{\X}\bigr).
\]
Inserting the same bounds $[\phi]_{\X}\le C$, $[\psi]_{\X}\le C(1+\lambda_n U)$ and
using $\lambda_n^{-9/8}\lambda_n=\lambda_n^{-1/8}\le1$ gives
\[
\|I^{\lambda_n}_{s,t}(x)-I^{\lambda_n}_{s,t}(z)\|_{L_m(\Omega)}
\le C|x-z|^{\mu}\bigl(U+\lambda_n^{-\frac18}U+\lambda_n^{-\frac98}\bigr)
\le C|x-z|^{\mu}\bigl(\lambda_n^{-\frac98}+U\bigr),
\]
which is \eqref{sdeb3} (note $\lambda_n^{-1/8}U\le U$). This completes the proof.
\end{proof}


\begin{thebibliography}{BDG25}

\bibitem[BDG25]{BDGLevy}
Oleg Butkovsky, Konstantinos Dareiotis, and M\'at\'e Gerencs\'er.
\newblock Strong rate of convergence of the {E}uler scheme for {SDE}s with
  irregular drift driven by {L}\'evy noise.
\newblock {\em Ann. Inst. Henri Poincar\'e{} Probab. Stat.}, 61(4), 2025.

\bibitem[BFG21]{BFG}
Carlo Bellingeri, Peter~K Friz, and M{\'a}t{\'e} Gerencs{\'e}r.
\newblock Singular paths spaces and applications.
\newblock {\em Stochastic Analysis and Applications}, pages 1--24, 2021.

\bibitem[BKS20]{BKS18}
Oleg Butkovsky, Alexei Kulik, and Michael Scheutzow.
\newblock Generalized couplings and ergodic rates for {SPDE}s and other
  {M}arkov models.
\newblock {\em Ann. Appl. Probab.}, 30(1):1--39, 2020.

\bibitem[BW20]{BW20}
Oleg Butkovsky and Fabrice Wunderlich.
\newblock Asymptotic strong feller property and local weak irreducibility via
  generalized couplings.
\newblock {\em arXiv preprint arXiv:1912.06121}, 2020.

\bibitem[HM06]{HM06}
Martin Hairer and Jonathan~C. Mattingly.
\newblock Ergodicity of the 2{D} {N}avier-{S}tokes equations with degenerate
  stochastic forcing.
\newblock {\em Ann. of Math. (2)}, 164(3):993--1032, 2006.

\bibitem[Shi96]{ShiryaevProb}
Albert~N. Shiryaev.
\newblock {\em Probability}.
\newblock Graduate Texts in Mathematics 95, Springer, 2nd edition, 1996.

\bibitem[L{\^e}20]{LeSSL}
Khoa L{\^e}.
\newblock A stochastic sewing lemma and applications.
\newblock {\em Electron. J. Probab.}, 25:Paper No. 38, 55, 2020.

\end{thebibliography}
\end{document}

